% s07_attractor.tex -- the global attractor is a convex planar graph; exact planar reduction.
% Exported labels: sec:attractor thm:attractor thm:twoorbit lem:absorbing prop:riccatientire prop:crestaxis
% cor:criterion.  Internal labels carry the prefix att:.
\section{The global attractor is a convex planar graph}\label{sec:attractor}
Theorem~\ref{thm:support} compares a cycle or the equilibrium with another log-bounded entire solution. This section extends the comparison to an arbitrary pair of log-bounded entire orbits and derives its geometric consequences. For all $\mu,k>0$ the union $\cA$ of these orbits is a compact global attractor; the projection $\pi(U,p,y)=(U,p)$ is injective on it; $\cA$ is the graph over $\Dsc=\pi(\cA)$ of the restriction $h$ of a convex function $H$; and the flow on $\cA$ is conjugate to the planar flow of $F(u,v)=(v,k-\mu v-H(u,v))$, whose divergence is $\eta-\mu$ on the interior of $\Dsc$. For $k>\kH$ the set $\Dsc$ is the closed Jordan domain bounded by the outermost projected cycle, and the cycle is unique at every parameter at which all cycles have $\lambda_c<0$ (Corollary~\ref{cor:criterion}). The sign of $\lambda_c$ is not determined in this section, and uniqueness of the cycle is not proved here.

Parts (A1)--(A5) of Theorem~\ref{thm:attractor} hold for all positive parameters. The proofs of (A1)--(A5) import only the Birkhoff--Hopf contraction theorem, which is (I1) in the list of Section~\ref{att:sec:setting}. They do not use Smith's permanence theorem \cite{Smith12}: dissipativity and permanence are proved directly, with explicit constants, in Lemma~\ref{lem:absorbing}. Parts (A6)--(A8) assume $k>\kH$, and their proofs use in addition the classical results (I2)--(I5) of that list: planar Poincar\'e--Bendixson theory, invariant manifolds and orbital stability, the topology of the plane, and Whitney's extension theorem. The competitive Poincar\'e--Bendixson theorem and the finiteness of the set of cycles \cite{Smith12} are not used in this section.

\paragraph{Local symbols.} Points of $\R^3$ are written in bold, $\mathbf a,\mathbf b,\mathbf z$, with coordinates $(U_{\mathbf a},p_{\mathbf a},y_{\mathbf a})$; they are unrelated to the components of $e_s$. Tangent vectors of $\R^3$ and solutions of the variational equation are written $\mathbf v=(v_1,v_2,v_3)$. $\mathbf f$ is the vector field of \eqref{eq:logsystem} and $\varphi_t$ its flow. Points of the $(U,p)$-plane are $q=(u,v)$, and $\chi_t$ is the flow of the reduced field $F$. A periodic orbit of $\chi_t$ in $\Dsc$ is denoted $\mathcal J$, its crest abscissa $\xi$, the set of crests $\mathrm{Cr}$; $\operatorname{ins}\mathcal J$ and $\operatorname{ext}\mathcal J$ are the bounded and the unbounded component of $\R^2\setminus\mathcal J$, and $\Omega=\operatorname{ins}\mathcal J_{\rm out}$ for the outermost projected cycle. $g=(\zeta,-\eta)\circ(\pi|_{\cA})^{-1}$ is the gradient jet of $h$, with modulus of continuity $\omega_g$. On the crest half-axis, $g_\ast(u)=h(u,0)-k$, $G_\ast$ is its primitive, $R_1$ and $T_1$ are the first-return map and time, and $\iota$ is the divergence integral over one turn. $\Phi(t,s)$ is the evolution operator of $\mathbf H'=B\mathbf H$, $d$ is Hilbert's projective metric, and $\underline m,\overline m,d_0,\nu_0$ are the constants \eqref{att:eq:birkhoff}. $x_\pm$ are the extreme values of $x=ke^{U}$ on $\cA$. $\mathbf M$ is the $3\times3$ monodromy matrix of a cycle; the matrix $M$ of Proposition~\ref{prop:crestmap} is its crest-section counterpart. $J$, $B$, $n$ and $O$ keep the meanings of Section~\ref{sec:adjoint}; here $n$ and $O$ are defined on all of $\cA$. Integers used as indices are $i,j$.

\subsection{Setting and imported classical results}\label{att:sec:setting}
Fix $\mu>0$ and $k>0$. Let $\mathbf f$ be the vector field of \eqref{eq:logsystem} on $\R^3$ and $\varphi_t$ its flow. In the coordinates $(U,p,y)$ the chemical domain is $\Dom=\{y>0,\ p<\lambda\}$. A \emph{log-bounded entire orbit} is the image of a solution of \eqref{eq:logsystem} that is defined and bounded on all of $\R$; positivity is not assumed. Every cycle and the equilibrium $E$ are log-bounded entire orbits.

For a positive continuous function $x(t)$, the \emph{Riccati system driven by $x$} is
\begin{equation}\label{att:eq:ric}
\eta'=\zeta+(\mu-1)\eta-\eta^{2},\qquad \zeta'=x-(1+\eta)\zeta ;
\end{equation}
along a solution of \eqref{eq:logsystem}, with $x=ke^{U}$, it is \eqref{eq:projective}. For a pair $(\eta,\zeta)$ put $n=(\zeta,-\eta,-1)^{\mathsf T}$. With $J$ as in \eqref{eq:adjointmatrices}, $J^{\mathsf T}n=(-x,\ \zeta+\mu\eta,\ \eta+1)^{\mathsf T}$, so the three components of $n'+J^{\mathsf T}n+(1+\eta)n$ are $\zeta'-x+(1+\eta)\zeta$, $-\eta'+\zeta+(\mu-1)\eta-\eta^2$ and $0$. Hence \eqref{att:eq:ric} is equivalent to
\begin{equation}\label{att:eq:normal}
n'+J^{\mathsf T}n+(1+\eta)n=0 .
\end{equation}

The following classical results are imported. The check that each hypothesis holds is made at the point of use.
\begin{enumerate}[label=\textup{(I\arabic*)}]
\item \emph{Birkhoff--Hopf theorem} \cite{Birkhoff57,Bushell73}; see also \cite[Chapter~3]{Seneta81}. Let $d(\mathbf u,\mathbf w)=\log\bigl(\max_i\frac{u_i}{w_i}\,\max_i\frac{w_i}{u_i}\bigr)$ be Hilbert's projective metric on $\inter\R^3_+$. An entrywise positive matrix $\mathbf Q$ contracts $d$ by the factor $\tanh(\Delta(\mathbf Q)/4)$, where $\Delta(\mathbf Q)=\sup\{d(\mathbf Q\mathbf u,\mathbf Q\mathbf w):\mathbf u,\mathbf w\in\R^3_+\setminus\{0\}\}$.
\item \emph{Planar theory} for $q'=F(q)$ with $F$ locally Lipschitz on $\R^2$ \cite[Chapter~VII]{Hartman02}.
\begin{enumerate}[label=\textup{(\alph*)}]
\item (Poincar\'e--Bendixson theorem with finitely many equilibria.) If a forward orbit is bounded and its $\omega$-limit set $\omega(q)$ contains finitely many equilibria, then $\omega(q)$ is an equilibrium, or a periodic orbit, or consists of equilibria together with orbits each of which converges to an equilibrium as $t\to+\infty$ and as $t\to-\infty$.
\item (Monotonicity on transversals.) A \emph{transversal} of $F$ is a closed line segment at no point of which $F$ vanishes or is parallel to the segment. By continuity $F$ then points to the same side of the segment at all of its points, so every orbit crosses a transversal in the same direction. Successive intersections of an orbit with a transversal are monotone along it: if an orbit meets a transversal at times $t_1<t_2<t_3$, then either the three points of intersection coincide and the orbit is periodic, or they are distinct and the second lies between the first and the third. Consequently a periodic orbit meets a transversal in at most one point. Neither statement holds for a segment that is merely transverse to the orbit at the points of intersection; a segment through an equilibrium of a rotation is an example. At each use below the segment is therefore shown to be a transversal in the sense just defined.
\end{enumerate}
\item \emph{Invariant manifolds and orbital stability} \cite[Chapter~IX]{Hartman02}.
\begin{enumerate}[label=\textup{(\alph*)}]
\item At a hyperbolic equilibrium of a $C^1$ vector field there are local stable and unstable manifolds, which are $C^1$ embedded discs tangent at the equilibrium to the stable and the unstable subspace. An orbit converging to the equilibrium as $t\to+\infty$ lies in the local stable manifold for all large $t$. Every orbit through a point of the local unstable manifold converges to the equilibrium as $t\to-\infty$.
\item The same statements hold at a hyperbolic fixed point of a $C^1$ local diffeomorphism; they are applied to the first-return map of a cycle on a transverse disc. For a cycle whose transverse multipliers satisfy $m_s<1<m_t$, flowing the local unstable curve of the return map gives the local unstable manifold of the cycle: a $C^1$ surface containing the cycle, whose tangent plane at a point $\mathbf a_0$ of the cycle is spanned by $\mathbf f(\mathbf a_0)$ and the eigenvector of the monodromy matrix for $m_t$, and whose points have backward orbits converging to the cycle.
\item A cycle all of whose transverse multipliers have modulus less than one is orbitally asymptotically stable: every neighbourhood $V$ of the cycle contains a neighbourhood $V'$ whose forward orbits remain in $V$ and have the cycle as $\omega$-limit set.
\end{enumerate}
\item \emph{Topology of the plane.} The Jordan curve theorem; the Schoenflies theorem, by which the closure of the bounded complementary component of a Jordan curve is the image of the closed unit disc under a homeomorphism that maps the unit circle onto the curve \cite{Thomassen92}; the fact that a Jordan curve, parametrised by a homeomorphism of the circle, has winding number $\pm1$ about each point of its bounded complementary component and $0$ about each point of the unbounded one; and invariance of the winding number about a point $q_0$ under homotopies of loops in $\R^2\setminus\{q_0\}$. In particular, a loop that extends to a continuous map of the closed unit disc into $\R^2\setminus\{q_0\}$ has winding number $0$ about $q_0$: the restrictions of this map to the circles of radius $r$ about the centre, $0\le r\le1$, form a homotopy in $\R^2\setminus\{q_0\}$ from the given loop to a constant loop. Invariance of domain is not needed: wherever a projected set is shown to be open, the inverse function theorem is used.
\item \emph{Whitney's extension theorem} in class $C^1$ \cite{Whitney34}: if $h$ and $g=(g_1,g_2)$ are continuous on a compact set $Q\subset\R^2$ and $h(q')-h(q)-g(q)\cdot(q'-q)=o(|q'-q|)$ uniformly on $Q$, then there is $\widehat h\in C^1(\R^2)$ with $\widehat h=h$ and $\nabla\widehat h=g$ on $Q$. It is used only in Theorem~\ref{thm:attractor}(A8)(a), at the outermost cycle.
\end{enumerate}
Existence, uniqueness and continuous dependence for ordinary differential equations, the implicit and inverse function theorems and Liouville's formula are used without further reference. From Sections~\ref{sec:adjoint}--\ref{sec:geometry} only Proposition~\ref{prop:multipliers} and Theorem~\ref{thm:geometry} enter, and both are re-derived below in the form needed (the first paragraph of Section~\ref{att:sec:mult} and Lemma~\ref{att:lem:cyclegeom}).

\subsection{Statement}\label{att:sec:statement}
\begin{theorem}[Attractor graph]\label{thm:attractor}
Let $\mu>0$, $k>0$.
\begin{enumerate}[label=\textup{(A\arabic*)}]
\item \textup{(Global attractor.)} $\Dom$ is forward invariant, and the semiflow on $\Dom$ has a compact, connected, invariant global attractor $\cA\subset\Dom$ which attracts every compact subset of $\Dom$. The set $\cA$ is exactly the union of all log-bounded entire orbits of \eqref{eq:logsystem} in $\R^3$.
\item \textup{(Positive adjoint ratios.)} Along every log-bounded entire orbit, with $x=ke^{U}$, the Riccati system \eqref{att:eq:ric} has exactly one solution on $\R$ with $\eta\ge0$ and $\zeta\ge0$. It is bounded,
\begin{equation}\label{att:eq:bounds}
\eta_-\le\eta\le\eta_+,\qquad \zeta_-\le\zeta\le\zeta_+,\qquad \eta_->\max(0,\mu-1),\quad \zeta_->0,
\end{equation}
with the explicit constants \eqref{att:eq:constants}, which depend only on $\mu$ and on $x_-=\min_{\cA}ke^{U}$, $x_+=\max_{\cA}ke^{U}$. It attracts all nonnegative data, in the pull-back and in the forward sense, at the uniform exponential rate \eqref{att:eq:pullback}. The resulting functions $\eta,\zeta:\cA\to(0,\infty)$, defined at $\mathbf a$ as the values at time $0$ for the orbit through $\mathbf a$, are continuous, and $t\mapsto(\eta,\zeta)(\varphi_t\mathbf a)$ is the Riccati solution along the orbit of $\mathbf a$. On a cycle they are the periodic ratios \eqref{eq:ratios}. At $E$ they are the constants $\eta_E=\gamma_E-1$, $\zeta_E=(\gamma_E-1)(\gamma_E-\mu)$, where $\gamma_E>\max(1,\mu)$ is the root of $\gamma_E(\gamma_E-1)(\gamma_E-\mu)=k$.
\item \textup{(Supporting planes.)} For all $\mathbf a\neq\mathbf b$ in $\cA$,
\begin{equation}\label{att:eq:support}
y_{\mathbf b}-y_{\mathbf a}>\zeta(\mathbf a)\,(U_{\mathbf b}-U_{\mathbf a})-\eta(\mathbf a)\,(p_{\mathbf b}-p_{\mathbf a}).
\end{equation}
\item \textup{(Convex $C^1$ graph.)} $\pi|_{\cA}$ is a homeomorphism onto the compact connected set $\Dsc=\pi(\cA)$. The function
\[
H(u,v)=\max_{\mathbf a\in\cA}\bigl\{y_{\mathbf a}+\zeta(\mathbf a)(u-U_{\mathbf a})-\eta(\mathbf a)(v-p_{\mathbf a})\bigr\}
\]
is convex and globally Lipschitz on $\R^2$, and $\cA=\{(u,v,h(u,v)):(u,v)\in\Dsc\}$ with $h=H|_{\Dsc}$. The map $g=(\zeta,-\eta)\circ(\pi|_\cA)^{-1}:\Dsc\to\R^2$ is continuous and, for all $q,q'\in\Dsc$,
\begin{equation}\label{att:eq:whitney}
0\le h(q')-h(q)-g(q)\cdot(q'-q)\le\bigl(g(q')-g(q)\bigr)\cdot(q'-q).
\end{equation}
Consequently $h$ is differentiable at every interior point of $\Dsc$, $h\in C^1(\inter\Dsc)$,
\[
\nabla h=(\zeta,-\eta),\qquad h_u=\zeta>0,\qquad h_v=-\eta<-\max(0,\mu-1)\qquad\text{on }\inter\Dsc,
\]
and $(h,g)$ satisfies the hypothesis of \textup{(I5)} on $\Dsc$. The set $\Dsc$ is not claimed to be convex: the graph is convex in the sense that $h$ is the restriction of the convex function $H$.
\item \textup{(Reduced planar flow.)} The field $F(u,v)=(v,\ k-\mu v-H(u,v))$ is globally Lipschitz on $\R^2$. Its flow $\chi_t$ leaves $\Dsc$ invariant, and $\pi\circ\varphi_t=\chi_t\circ\pi$ on $\cA$. Every $(u,v)\in\Dsc$ satisfies
\begin{equation}\label{att:eq:pde}
\zeta\,v-\eta\,(k-\mu v-h)=ke^{u}-h ,
\end{equation}
that is, $h_uv+h_v(k-\mu v-h)=ke^{u}-h$ on $\inter\Dsc$. On $\inter\Dsc$ the field $F$ is $C^1$ with $\diver F=\eta-\mu$. Moreover $K=\zeta U-\eta p-y+k>0$ and $W=\zeta'p-\eta'p'>0$ on $\cA\setminus\{E\}$, where the prime is the derivative along the flow, and $K'+(1+\eta)K=k\,\psiU(U)$, $W'+(1+\eta)W=xp^2$ on $\cA$.
\end{enumerate}
Assume in addition $k>\kH$. Then:
\begin{enumerate}[label=\textup{(A\arabic*)},resume]
\item \textup{($E$ is a source in $\cA$.)} With the cones $C_\pm$ of \eqref{eq:cones}, $\cA\cap(E+C_\pm)=\{E\}$. No log-bounded entire orbit other than $E$ converges to $E$ as $t\to+\infty$. The two-dimensional unstable manifold $W^u(E)$ lies in $\cA$, and $(0,0)\in\inter\Dsc$. Every orbit in $\cA\setminus\{E\}$ has a cycle as $\omega$-limit set; its $\alpha$-limit set is $\{E\}$ or a cycle; and no non-periodic orbit has the same cycle as $\alpha$- and $\omega$-limit set. In particular cycles exist.
\item \textup{(The domain.)} The set of crests $\mathrm{Cr}=\{u>0:(u,0)\text{ lies on a projected cycle}\}$ is a nonempty compact subset of $(0,\infty)$. Let $\Gamma_{\rm out}$ and $\Gamma_{\rm in}$ be the cycles with the largest and the smallest crest. Then $\Dsc$ is the closed Jordan domain bounded by $\pi(\Gamma_{\rm out})$: $\partial\Dsc=\pi(\Gamma_{\rm out})$, $\inter\Dsc$ is the bounded component of $\R^2\setminus\pi(\Gamma_{\rm out})$, and $\cA$ is a closed topological disc with boundary $\Gamma_{\rm out}$.
\item \textup{(Multipliers.)} Let $\Gamma$ be a cycle with period $T$, $m_t=\exp\int_0^T(\eta-\mu)\,\mathrm dt$ and $\lambda_c=\avg\eta-\mu$, as in Proposition~\ref{prop:multipliers}.
\begin{enumerate}[label=\textup{(\alph*)}]
\item If $\pi(\Gamma)\subset\inter\Dsc$, which holds for every cycle except $\Gamma_{\rm out}$, then $\pi(\Gamma)$ is a periodic orbit of the $C^1$ planar field $F$ and its planar Floquet multiplier $\exp\oint\diver F\,\mathrm dt$ equals $m_t$. For $\Gamma_{\rm out}$ the same holds for the $C^1$ field $\widehat F=(v,k-\mu v-\widehat h)$, where $\widehat h$ is any extension of the jet $(h,g)$ given by \textup{(I5)}; $\widehat F=F$ on $\Dsc$.
\item If $m_t>1$, no orbit in $\cA\setminus\Gamma$ has $\omega$-limit set $\Gamma$. If $m_t<1$, no orbit in $\cA\setminus\Gamma$ has $\alpha$-limit set $\Gamma$.
\item $\lambda_c(\Gamma_{\rm in})\le0$ and $\lambda_c(\Gamma_{\rm out})\le0$.
\item If every cycle has $\lambda_c<0$, there is exactly one cycle. Equivalently, if there are at least two cycles, some cycle has $\lambda_c\ge0$.
\item If every cycle is hyperbolic, that is $\lambda_c\neq0$, the number of cycles is finite and odd, and ordered by nesting they alternate $\lambda_c<0$, $\lambda_c>0$, \dots, $\lambda_c<0$.
\item \textup{(Dulac form.)} Let $\Lambda$ be a real function on the union of all cycles such that $t\mapsto\Lambda(\varphi_t\mathbf a)$ is absolutely continuous, with derivative $\Dt\Lambda$. If $\Dt\Lambda+\eta-\mu\le0$ almost everywhere on every cycle and is not almost everywhere zero on any cycle, the cycle is unique and hyperbolic attracting. There is no $\Lambda$ for which $\Dt\Lambda+\eta-\mu\ge0$ almost everywhere on $\Gamma_{\rm in}$, or almost everywhere on $\Gamma_{\rm out}$, without vanishing almost everywhere there. If $\Dt\Lambda+\eta-\mu=0$ almost everywhere on a cycle, that cycle has $\lambda_c=0$; this case is not excluded.
\end{enumerate}
\end{enumerate}
\end{theorem}
For $0<k\le\kH$ the equilibrium is globally asymptotically stable in $\Dom$ by \cite[Theorem~1.1(b)]{Smith12} and, at $k=\kH$, by \cite[Theorem~13]{BH23}. A globally asymptotically stable equilibrium attracts compact sets, and $\cA=\varphi_t(\cA)$, so $\cA=\{E\}$ and (A1)--(A5) reduce to statements about $E$. This remark is imported and is not used below.

\subsection{Dissipativity and the global attractor: proof of (A1)}\label{att:sec:diss}
We use $x=ke^{U}>0$, $\ell=\lambda-p$ and $y$. Since $\mu\lambda=k$, \eqref{eq:logsystem} reads
\begin{equation}\label{att:eq:xyl}
x'=(\lambda-\ell)x,\qquad y'=x-y,\qquad \ell'=y-\mu\ell,
\end{equation}
and $\Dom=\{x>0,\ y>0,\ \ell>0\}$.

\begin{lemma}[Explicit absorbing box]\label{lem:absorbing}
Every solution starting in $\Dom$ exists for all $t\ge0$ and stays in $\Dom$. There are constants $0<\underline x<\overline x$, $0<\underline y<\overline y$, $0<\underline\ell<\overline\ell$, depending only on $(\mu,k)$, such that every solution in $\Dom$ satisfies, for all sufficiently large $t$,
\[
(x,y,\ell)(t)\in\mathcal B_0:=[\underline x,\overline x]\times[\underline y,\overline y]\times[\underline\ell,\overline\ell].
\]
\end{lemma}
\begin{proof}
On the maximal interval of existence $x(t)=x(0)\exp\int_0^t(\lambda-\ell)\,\mathrm dr>0$ and, by variation of constants, for $s\le t$,
\begin{equation}\label{att:eq:voc}
y(t)=e^{-(t-s)}y(s)+\int_s^te^{-(t-r)}x(r)\,\mathrm dr,\qquad \ell(t)=e^{-\mu(t-s)}\ell(s)+\int_s^te^{-\mu(t-r)}y(r)\,\mathrm dr ,
\end{equation}
so $y>0$ and $\ell>0$ persist. As long as $\ell>0$ we have $x'\le\lambda x$, so $x\le x(0)e^{\lambda t}$, and then $y$ and $\ell$ are bounded on bounded intervals by \eqref{att:eq:voc}; in turn $x\ge x(0)\exp(-\int_0^t\ell\,\mathrm dr)$ is bounded away from $0$ on bounded intervals. Hence solutions are global in forward time and $\Dom$ is forward invariant.

\emph{Upper bounds.} From $x'/x\le\lambda$ we get $x(r)\ge x(t)e^{-\lambda(t-r)}$ for $0\le r\le t$. Dropping the nonnegative initial terms in \eqref{att:eq:voc} with $s=0$, twice,
\begin{gather*}
\ell(t)\ge\int_0^te^{-\mu(t-s)}\!\!\int_0^se^{-(s-r)}x(r)\,\mathrm dr\,\mathrm ds\ \ge\ \beta(t)\,x(t),\\
\beta(t):=\iint_{r_1,r_2\ge0,\ r_1+r_2\le t}e^{-\mu r_1}e^{-r_2}e^{-\lambda(r_1+r_2)}\,\mathrm dr_1\,\mathrm dr_2 ,
\end{gather*}
with $r_1=t-s$ and $r_2=s-r$. The function $\beta$ is positive and nondecreasing; put $\beta_1=\beta(1)$. For $t\ge1$, $x'\le x(\lambda-\beta_1x)$, and comparison with the logistic equation gives $\limsup x\le\lambda/\beta_1$. Put $\overline x=2\lambda/\beta_1$; then $x\le\overline x$ for $t\ge t_1$. By \eqref{att:eq:voc}, $y(t)\le e^{-(t-t_1)}y(t_1)+\overline x$, so $y\le\overline y:=2\overline x$ for $t\ge t_2$, and likewise $\ell\le\overline\ell:=2\overline y/\mu$ for $t\ge t_3$.

\emph{Lower bounds.} Let $\mu_0=\min(1,\mu)$ and choose $L>0$ with $(\overline\ell+\overline ys)e^{-\mu_0s}\le\lambda/4$ for all $s\ge L$. Put $\underline x=\tfrac k4\,e^{-\overline\ell L}$. Consider $t_0\ge t_3$ and an interval $[t_0,t_0+\vartheta]$, $0<\vartheta\le\infty$, on which $x\le k/4$. For $0\le s\le\vartheta$, \eqref{att:eq:voc} gives $y(t_0+s)\le\overline ye^{-s}+k/4$ and then
\[
\ell(t_0+s)\le\overline\ell e^{-\mu s}+\overline y\int_0^se^{-\mu(s-r)}e^{-r}\,\mathrm dr+\frac{k}{4\mu}\le(\overline\ell+\overline ys)e^{-\mu_0s}+\frac\lambda4 ,
\]
because $\int_0^se^{-\mu(s-r)}e^{-r}\,\mathrm dr\le\int_0^se^{-\mu_0(s-r)}e^{-\mu_0r}\,\mathrm dr=se^{-\mu_0s}$. Hence $\ell\le\lambda/2$ and $x'\ge(\lambda/2)x$ on $[t_0+L,t_0+\vartheta]$. Consequently $\vartheta<\infty$, since otherwise $x\to\infty$, contradicting $x\le k/4$: every maximal interval in $[t_3,\infty)$ on which $x\le k/4$ is bounded. Let $t_4\ge t_3$ be a time with $x(t_4)\ge k/4$, and let $[t_0,t_0+\vartheta]\subset[t_4,\infty)$ be a maximal interval with $x\le k/4$; then $x(t_0)=k/4$. On $[t_0,t_0+\min(\vartheta,L)]$, $x'/x=\lambda-\ell\ge-\overline\ell$ gives $x\ge\tfrac k4e^{-\overline\ell L}=\underline x$; on the rest of the interval $x$ is nondecreasing. So $x\ge\underline x$ on $[t_4,\infty)$. Then \eqref{att:eq:voc} gives $y(t)\ge\underline x(1-e^{-(t-t_4)})\ge\underline y:=\underline x/2$ eventually, and then $\ell\ge\underline\ell:=\underline y/(2\mu)$ eventually. The times $t_1,\dots,t_4$ depend on the solution; the constants do not. They are ordered as stated: $\beta_1<\int_0^1e^{-\mu r_1}\,\mathrm dr_1<1/\mu$, so $\overline x>2\lambda\mu=2k>k/4>\underline x$, and then $\underline y=\underline x/2<2\overline x=\overline y$ and $\underline\ell=\underline y/(2\mu)<2\overline y/\mu=\overline\ell$.
\end{proof}
In the coordinates $(U,p,y)$, $\mathcal B_0$ is the compact box $[\log(\underline x/k),\log(\overline x/k)]\times[\lambda-\overline\ell,\lambda-\underline\ell]\times[\underline y,\overline y]\subset\Dom$.

\begin{proof}[Proof of (A1)]
Let $\mathcal B'\subset\Dom$ be a compact box, in the coordinates $(U,p,y)$, whose interior contains $\mathcal B_0$.

\emph{(a) A compact forward-invariant absorbing set.} For $\mathbf z\in\Dom$ Lemma~\ref{lem:absorbing} gives $t_{\mathbf z}\ge1$ with $\varphi_{t_{\mathbf z}}(\mathbf z)\in\mathcal B_0\subset\inter\mathcal B'$ and, by continuity, an open set $V_{\mathbf z}\ni\mathbf z$ in $\Dom$ with $\varphi_{t_{\mathbf z}}(V_{\mathbf z})\subset\inter\mathcal B'$. Cover $\mathcal B'$ by finitely many $V_{\mathbf z_i}$, let $t_\ast=\max_it_{\mathbf z_i}$, and put $\mathcal B_1=\varphi([0,t_\ast]\times\mathcal B')$, a compact subset of $\Dom$ which contains $\mathcal B'$.

We first show that $\varphi_t(\mathbf z)\in\mathcal B_1$ for all $\mathbf z\in\mathcal B'$ and all $t\ge0$. Put $s_0=0$. If $s_j$ has been defined and $\varphi_{s_j}(\mathbf z)\in\mathcal B'$, choose $i$ with $\varphi_{s_j}(\mathbf z)\in V_{\mathbf z_i}$ and put $s_{j+1}=s_j+t_{\mathbf z_i}$. Then $\varphi_{s_{j+1}}(\mathbf z)\in\mathcal B'$, and for $s_j\le t\le s_{j+1}$ the point $\varphi_t(\mathbf z)=\varphi_{t-s_j}(\varphi_{s_j}(\mathbf z))$ lies in $\varphi([0,t_\ast]\times\mathcal B')=\mathcal B_1$, because $0\le t-s_j\le t_{\mathbf z_i}\le t_\ast$. By induction this defines $s_j$ for all $j$, and $s_{j+1}-s_j\ge1$, so the intervals $[s_j,s_{j+1}]$ cover $[0,\infty)$.

$\mathcal B_1$ is forward invariant: a point of $\mathcal B_1$ has the form $\varphi_r(\mathbf b)$ with $\mathbf b\in\mathcal B'$ and $0\le r\le t_\ast$, and $\varphi_t(\varphi_r(\mathbf b))=\varphi_{t+r}(\mathbf b)\in\mathcal B_1$ for all $t\ge0$ by the previous paragraph. $\mathcal B_1$ absorbs compact sets: a compact $Q\subset\Dom$ is covered by finitely many of the sets $V_{\mathbf z}$, $\mathbf z\in Q$; let $t_Q$ be the largest of the corresponding times $t_{\mathbf z}$. A point $\mathbf w\in Q$ lies in one of these sets $V_{\mathbf z}$, so $\varphi_{t_{\mathbf z}}(\mathbf w)\in\mathcal B'$ and, by the previous paragraph, $\varphi_t(\mathbf w)\in\mathcal B_1$ for all $t\ge t_{\mathbf z}$, in particular for all $t\ge t_Q$. Hence $\varphi_t(Q)\subset\mathcal B_1$ for all $t\ge t_Q$.

\emph{(b) The attractor.} Since $\mathcal B_1$ is forward invariant, $\varphi_{t'}(\mathcal B_1)=\varphi_t(\varphi_{t'-t}(\mathcal B_1))\subset\varphi_t(\mathcal B_1)$ for $t'\ge t\ge0$. Put $\cA=\bigcap_{t\ge0}\varphi_t(\mathcal B_1)$, a decreasing intersection of nonempty compact sets. Then $\operatorname{dist}(\varphi_t(\mathcal B_1),\cA)\to0$, so $\cA$ attracts $\mathcal B_1$ and, by (a), every compact subset of $\Dom$. Invariance: if $\mathbf z\in\cA$ and $s\ge0$ then $\varphi_s(\mathbf z)\in\varphi_{s+t}(\mathcal B_1)$ for all $t\ge0$ and $\varphi_s(\mathbf z)\in\varphi_s(\mathcal B_1)\subset\varphi_{t'}(\mathcal B_1)$ for $t'\le s$, so $\varphi_s(\cA)\subset\cA$. Conversely, for $\mathbf z\in\cA$ and integers $j\ge s$ write $\mathbf z=\varphi_j(\mathbf b_j)=\varphi_s(\varphi_{j-s}(\mathbf b_j))$ with $\mathbf b_j\in\mathcal B_1$; by injectivity of $\varphi_s$ the point $\mathbf w=\varphi_{j-s}(\mathbf b_j)$ does not depend on $j$, and it lies in $\varphi_{j-s}(\mathcal B_1)$ for all integers $j\ge s$. Given $t\ge0$, choose such a $j$ with $j-s\ge t$; since the sets decrease, $\mathbf w\in\varphi_{j-s}(\mathcal B_1)\subset\varphi_t(\mathcal B_1)$. Hence $\mathbf w\in\cA$, and $\varphi_s(\mathbf w)=\mathbf z$. So $\varphi_s(\cA)=\cA$, and each point of $\cA$ lies on an entire orbit contained in the compact set $\cA$.

\emph{(c) Characterisation.} Let $\mathbf z(\cdot)=(U,p,y)(\cdot)$ be a solution of \eqref{eq:logsystem} on $\R$ that is bounded. Then $x=ke^{U}$ satisfies $x_{\inf}:=\inf_tx(t)>0$. The bounded function $y$ solves $y'=x-y$ on $\R$, and two bounded solutions of this equation differ by a multiple of $e^{-t}$, which is bounded on $\R$ only if it vanishes; hence $y(t)=\int_{-\infty}^te^{-(t-s)}x(s)\,\mathrm ds\ge x_{\inf}$. Likewise $\ell=\lambda-p=\int_{-\infty}^te^{-\mu(t-s)}y(s)\,\mathrm ds\ge x_{\inf}/\mu$. So the closure $Q$ of the orbit is a compact subset of $\Dom$, and for every $s$ and every $t\ge0$, $\mathbf z(s)=\varphi_t\bigl(\varphi_{t_Q}(\mathbf z(s-t-t_Q))\bigr)\in\varphi_t(\mathcal B_1)$. Thus $\mathbf z(s)\in\cA$. With (b) this proves the characterisation.

\emph{(d) Connectedness.} Let $\mathcal B''\subset\Dom$ be a compact box containing $\mathcal B_1$. Then $\varphi_t(\mathcal B'')$ is connected, contains $\cA=\varphi_t(\cA)$, and lies in any prescribed neighbourhood of $\cA$ for large $t$. If $\cA$ were the union of two disjoint nonempty compact sets, separated by disjoint open sets $O_1,O_2$, then for large $t$ the connected set $\varphi_t(\mathcal B'')\subset O_1\cup O_2$ would meet both, a contradiction.
\end{proof}

\subsection{The positive adjoint direction: proof of (A2)}\label{att:sec:adjoint}
Throughout this subsection $x:\R\to[x_-,x_+]$ is a continuous function with $0<x_-\le x_+$; no dynamics is assumed. Consider
\begin{equation}\label{att:eq:lin}
\mathbf H'=B(t)\mathbf H,\qquad B(t)=\begin{pmatrix}0&0&x(t)\\1&\mu&0\\0&1&1\end{pmatrix},
\end{equation}
which is the adjoint system of \eqref{eq:adjointmatrices} with a general driving function, and let $\Phi(t,s)$ be its evolution operator. If $\mathbf H$ solves \eqref{att:eq:lin} with $H_3>0$, then $\eta=H_2/H_3$, $\zeta=H_1/H_3$ solve \eqref{att:eq:ric}. Conversely, if $(\eta,\zeta)$ solves \eqref{att:eq:ric} on an interval, then $H_3=\exp\int(1+\eta)\,\mathrm dt$, $H_2=\eta H_3$, $H_1=\zeta H_3$ solves \eqref{att:eq:lin}; both statements follow by direct differentiation. Let $B_\pm$ be $B$ with $x$ replaced by $x_\pm$, and
\begin{equation}\label{att:eq:birkhoff}
\underline m=\min_{i,j}\bigl(e^{B_-}\bigr)_{ij},\qquad \overline m=\max_{i,j}\bigl(e^{B_+}\bigr)_{ij},\qquad
d_0=2\log\frac{\overline m}{\underline m},\qquad \nu_0=\tanh\frac{d_0}{4}=\frac{\overline m-\underline m}{\overline m+\underline m}<1 .
\end{equation}

\begin{lemma}[Uniform positivity and contraction]\label{att:lem:cone}
\begin{enumerate}[label=\textup{(\roman*)}]
\item $\underline m>0$, and for every $t$ all entries of $\Phi(t+1,t)$ lie in $[\underline m,\overline m]$. For $t\ge s$, $\Phi(t,s)$ is entrywise nonnegative.
\item For $t-s\ge1$ and $\mathbf u,\mathbf w\in\R^3_+\setminus\{0\}$,
$d\bigl(\Phi(t,s)\mathbf u,\Phi(t,s)\mathbf w\bigr)\le d_0\,\nu_0^{\lfloor t-s\rfloor-1}$.
\end{enumerate}
\end{lemma}
\begin{proof}
(i) All entries of $B(t)$ and $B_\pm$ are nonnegative (the diagonal is $0,\mu,1$), and $B_-\le B(t)\le B_+$ entrywise. The Peano--Baker series $\Phi(t,s)=\mathrm{Id}+\int_s^tB(t_1)\,\mathrm dt_1+\int_s^tB(t_1)\int_s^{t_1}B(t_2)\,\mathrm dt_2\,\mathrm dt_1+\cdots$ has nonnegative terms and is entrywise monotone in $B$, so $e^{(t-s)B_-}\le\Phi(t,s)\le e^{(t-s)B_+}$. The directed graph of $B_-$ contains the cycle $1\to2\to3\to1$, because the entries $(2,1)$, $(3,2)$, $(1,3)$ are positive; so $B_-$ is irreducible and $e^{B_-}=\sum_{j\ge0}B_-^j/j!$ is entrywise positive.

(ii) Let $\mathbf Q$ be a matrix with entries in $[\underline m,\overline m]$ and $\mathbf u,\mathbf w\in\R^3_+\setminus\{0\}$. Then
$(\mathbf Q\mathbf u)_i/(\mathbf Q\mathbf w)_i\in\bigl[\tfrac{\underline m|\mathbf u|_1}{\overline m|\mathbf w|_1},\tfrac{\overline m|\mathbf u|_1}{\underline m|\mathbf w|_1}\bigr]$ for each $i$, hence $d(\mathbf Q\mathbf u,\mathbf Q\mathbf w)\le d_0$ and $\Delta(\mathbf Q)\le d_0$; by (I1), $\mathbf Q$ contracts $d$ by the factor $\nu_0$. Put $j=\lfloor t-s\rfloor\ge1$ and factor
\[
\Phi(t,s)=\Phi(t,t-1)\cdots\Phi(t-j+2,t-j+1)\;\Phi(t-j+1,t-j)\;\Phi(t-j,s).
\]
The last factor maps $\R^3_+\setminus\{0\}$ into itself, being nonnegative and invertible; the next one maps it into $\inter\R^3_+$ with an image of $d$-diameter at most $d_0$; and each of the remaining $j-1$ factors contracts $d$ by $\nu_0$.
\end{proof}

\begin{lemma}[Suprema of bounded $C^1$ functions]\label{att:lem:sup}
If $\phi\in C^1(\R)$ is bounded, there are $t_j$ with $\phi(t_j)\to\sup\phi$ and $\phi'(t_j)\to0$, and similarly for $\inf\phi$.
\end{lemma}
\begin{proof}
Let $\overline\phi=\sup\phi$, pick $s_j$ with $\phi(s_j)\ge\overline\phi-j^{-2}$, and let $t_j$ be a maximum point of $\phi(t)-(t-s_j)^2/j$. Then $\phi(t_j)-(t_j-s_j)^2/j\ge\phi(s_j)\ge\overline\phi-j^{-2}$, whence $\phi(t_j)\ge\overline\phi-j^{-2}$, $(t_j-s_j)^2\le j^{-1}$ and $|\phi'(t_j)|=2|t_j-s_j|/j\le2j^{-3/2}$.
\end{proof}

\begin{proposition}[Unique positive entire Riccati solution]\label{prop:riccatientire}
Let $x:\R\to[x_-,x_+]$ be continuous, $0<x_-\le x_+$.
\begin{enumerate}[label=\textup{(\roman*)}]
\item There is exactly one solution $(\eta,\zeta)$ of \eqref{att:eq:ric} on $\R$ with $\eta\ge0$ and $\zeta\ge0$ everywhere.
\item It satisfies \eqref{att:eq:bounds} with the constants of Lemma~\ref{lem:projectivebounds},
\begin{equation}\label{att:eq:constants}
\eta_+=\tfrac{\mu-1+\sqrt{(\mu-1)^2+4x_+}}2,\quad \zeta_-=\tfrac{x_-}{1+\eta_+},\quad
\eta_-=\tfrac{\mu-1+\sqrt{(\mu-1)^2+4\zeta_-}}2,\quad \zeta_+=\tfrac{x_+}{1+\eta_-}.
\end{equation}
\item For every $s\in\R$ and $(\eta^\circ,\zeta^\circ)\in[0,\infty)^2$ the solution $(\tilde\eta,\tilde\zeta)$ of \eqref{att:eq:ric} with data $(\eta^\circ,\zeta^\circ)$ at time $s$ exists on $[s,\infty)$, and for $t\ge s+1$
\begin{equation}\label{att:eq:pullback}
\bigl|\log\tilde\eta(t)-\log\eta(t)\bigr|\le d_0\,\nu_0^{\lfloor t-s\rfloor-1},\qquad
\bigl|\log\tilde\zeta(t)-\log\zeta(t)\bigr|\le d_0\,\nu_0^{\lfloor t-s\rfloor-1}.
\end{equation}
\end{enumerate}
\end{proposition}
\begin{proof}
\emph{Existence.} Let $\triangle=\{\mathbf u\in\R^3_+:u_1+u_2+u_3=1\}$ and, for $s<t$, $\triangle(t,s)=\{\Phi(t,s)\mathbf u/|\Phi(t,s)\mathbf u|_1:\mathbf u\in\triangle\}$, a compact subset of $\triangle$. Since $\Phi(t,s')=\Phi(t,s)\Phi(s,s')$ and $\Phi(s,s')$ maps $\R^3_+\setminus\{0\}$ into itself, $\triangle(t,s')\subset\triangle(t,s)$ for $s'<s$. For $t-s\ge1$, $\triangle(t,s)\subset\inter\triangle$ and its $d$-diameter is at most $d_0\nu_0^{\lfloor t-s\rfloor-1}$ by Lemma~\ref{att:lem:cone}. Hence $\bigcap_{s<t}\triangle(t,s)$ is nonempty, as an intersection of nested compact sets, and has $d$-diameter $0$; since $d$ is a metric on $\inter\triangle$, it is a single point $\widehat{\mathbf H}(t)\in\inter\triangle$. For $t'>t$ the vector $\Phi(t',t)\widehat{\mathbf H}(t)$, normalised, lies in $\triangle(t',s)$ for every $s<t$, hence in $\bigcap_{s<t}\triangle(t',s)=\bigcap_{s<t'}\triangle(t',s)=\{\widehat{\mathbf H}(t')\}$. Therefore $\mathbf H^\ast(t):=\Phi(t,0)\widehat{\mathbf H}(0)$ is, for every $t\in\R$, a positive multiple of $\widehat{\mathbf H}(t)$: it is a positive entire solution of \eqref{att:eq:lin}. Its ratios $(\eta,\zeta)$ solve \eqref{att:eq:ric} on $\R$ and are positive.

\emph{Uniqueness.} Let $(\tilde\eta,\tilde\zeta)$ be a solution on $\R$ with $\tilde\eta,\tilde\zeta\ge0$. The associated $\widetilde{\mathbf H}$ is a solution of \eqref{att:eq:lin} with values in $\R^3_+\setminus\{0\}$, so the normalisation of $\widetilde{\mathbf H}(t)=\Phi(t,s)\widetilde{\mathbf H}(s)$ lies in $\triangle(t,s)$ for all $s<t$; it therefore equals $\widehat{\mathbf H}(t)$. Hence $(\tilde\eta,\tilde\zeta)=(\eta,\zeta)$.

\emph{Bounds.} $\mathbf H^\ast(t)=\Phi(t,t-1)\mathbf c$ with $\mathbf c=\mathbf H^\ast(t-1)\in\R^3_+\setminus\{0\}$, so $\underline m|\mathbf c|_1\le H^\ast_i(t)\le\overline m|\mathbf c|_1$ and $\eta,\zeta\in[\underline m/\overline m,\overline m/\underline m]$: both are bounded and bounded away from $0$, and by \eqref{att:eq:ric} their derivatives are bounded. Write $\overline\phi=\sup\phi$, $\underline\phi=\inf\phi$. Lemma~\ref{att:lem:sup} applied to $\zeta$ gives $t_j$ with $x(t_j)-(1+\eta(t_j))\zeta(t_j)\to0$ and $\zeta(t_j)\to\overline\zeta$, so $\overline\zeta\le x_+/(1+\underline\eta)\le x_+$. Applied to $\eta$ at its supremum, it gives $\overline\eta(\overline\eta+1-\mu)=\lim\zeta(t_j)\le x_+$ with $\overline\eta>0$, so $\overline\eta\le\eta_+$. At the infimum of $\zeta$: $\underline\zeta=\lim x(t_j)/(1+\eta(t_j))\ge x_-/(1+\eta_+)=\zeta_-$. At the infimum of $\eta$: $\underline\eta(\underline\eta+1-\mu)=\lim\zeta(t_j)\ge\zeta_->0$ with $\underline\eta>0$; the quadratic $s(s+1-\mu)$ is nonpositive on $[0,\max(0,\mu-1)]$ and increasing beyond, so $\underline\eta\ge\eta_->\max(0,\mu-1)$. Finally $\overline\zeta\le x_+/(1+\eta_-)=\zeta_+$.

\emph{Attraction.} The solution of \eqref{att:eq:lin} with $\widetilde{\mathbf H}(s)=(\zeta^\circ,\eta^\circ,1)^{\mathsf T}$ is nonnegative, and $\widetilde H_3'=\widetilde H_2+\widetilde H_3\ge\widetilde H_3$ keeps $\widetilde H_3>0$. Its ratios are the solution of \eqref{att:eq:ric} with the given data, which is unique because \eqref{att:eq:ric} is polynomial in $(\eta,\zeta)$, and they are defined on $[s,\infty)$. By Lemma~\ref{att:lem:cone}(ii) with $\mathbf u=\widetilde{\mathbf H}(s)$, $\mathbf w=\mathbf H^\ast(s)$, $d(\widetilde{\mathbf H}(t),\mathbf H^\ast(t))\le d_0\nu_0^{\lfloor t-s\rfloor-1}$, and $|\log(u_i/u_3)-\log(w_i/w_3)|=|\log(u_i/w_i)-\log(u_3/w_3)|\le d(\mathbf u,\mathbf w)$ for vectors in $\inter\R^3_+$.
\end{proof}
The estimate \eqref{att:eq:pullback} is used below only qualitatively, to obtain continuity. Since $1-\nu_0=2\underline m/(\overline m+\underline m)$, the factor $\nu_0$ is close to $1$ when the range of $x$ is large, and \eqref{att:eq:pullback} is then far from the contraction that actually takes place.

\begin{proof}[Proof of (A2)]
By (A1) a log-bounded entire orbit lies in $\cA$. For $\mathbf a\in\cA$ put $x_{\mathbf a}(t)=k\exp U(\varphi_t\mathbf a)$, $t\in\R$; then $x_-\le x_{\mathbf a}\le x_+$ with $x_\pm$ the extreme values of $ke^{U}$ on the compact invariant set $\cA$. Proposition~\ref{prop:riccatientire} gives the unique nonnegative entire solution $(\eta_{\mathbf a},\zeta_{\mathbf a})(t)$, the bounds, and \eqref{att:eq:pullback} with constants uniform in $\mathbf a$. Define $(\eta,\zeta)(\mathbf a)=(\eta_{\mathbf a},\zeta_{\mathbf a})(0)$. Since $x_{\varphi_s\mathbf a}(t)=x_{\mathbf a}(t+s)$, uniqueness gives $(\eta,\zeta)(\varphi_s\mathbf a)=(\eta_{\mathbf a},\zeta_{\mathbf a})(s)$.

\emph{Continuity.} For $s\ge1$ let $(\eta^{[s]},\zeta^{[s]})(\mathbf a)$ be the value at time $0$ of the solution of \eqref{att:eq:ric}, driven by $x_{\mathbf a}$, with data $(1,1)$ at time $-s$; it is the pair of ratios of the vector $\Phi_{\mathbf a}(0,-s)(1,1,1)^{\mathsf T}$, where $\Phi_{\mathbf a}$ is the evolution operator for the driving function $x_{\mathbf a}$. The flow restricted to the compact invariant set $\cA$ is a continuous action of $\R$, so $\mathbf a\mapsto x_{\mathbf a}|_{[-s,0]}$ is continuous into $C([-s,0])$, and by continuous dependence of the linear system \eqref{att:eq:lin} on its coefficient, $(\eta^{[s]},\zeta^{[s]}):\cA\to(0,\infty)^2$ is continuous. By \eqref{att:eq:pullback},
\[
\sup_{\mathbf a\in\cA}\bigl|\log(\eta^{[s]},\zeta^{[s]})(\mathbf a)-\log(\eta,\zeta)(\mathbf a)\bigr|\le d_0\,\nu_0^{\lfloor s\rfloor-1}\to0\qquad(s\to\infty),
\]
the logarithm being taken componentwise. A uniform limit of continuous functions is continuous.

\emph{Cycles and $E$.} On a cycle of period $T$, $(\eta,\zeta)(\cdot+T)$ is again a nonnegative entire solution, so $(\eta,\zeta)$ is $T$-periodic; the periodic ratios \eqref{eq:ratios} are positive entire solutions and hence coincide with it. At $E$ the constant solutions of \eqref{att:eq:ric} with $x=k$ satisfy $\zeta=k/(1+\eta)$ and $\eta(\eta+1-\mu)=\zeta$, that is, $\gamma(\gamma-1)(\gamma-\mu)=k$ with $\gamma=1+\eta$; positivity of $\eta$ and $\zeta$ forces $\gamma>\max(1,\mu)$, where the cubic increases from $0$ to $\infty$. By uniqueness this constant solution is the positive entire solution along $E$.
\end{proof}

\subsection{Supporting planes: proof of (A3)}\label{att:sec:support}
\begin{theorem}[Two-orbit comparison]\label{thm:twoorbit}
Let $z_1=(U_1,p_1,y_1)$ and $z_2=(U_2,p_2,y_2)$ be log-bounded entire solutions of \eqref{eq:logsystem} with $z_1\not\equiv z_2$, and let $(\eta,\zeta)$ be the positive entire Riccati solution along $z_1$, $n=(\zeta,-\eta,-1)^{\mathsf T}$, $x_1=ke^{U_1}$. Then
$\Delta:=n\cdot(z_2-z_1)=\zeta(U_2-U_1)-\eta(p_2-p_1)-(y_2-y_1)$ satisfies
\begin{equation}\label{att:eq:deltaode}
\Delta'+(1+\eta)\Delta=-x_1\,\Rem(U_2-U_1)\le0
\end{equation}
and $\Delta(t)<0$ for all $t\in\R$. No positivity, periodicity or stability of $z_1$ or $z_2$ is assumed, and $z_2$ may be a time translate of $z_1$.
\end{theorem}
\begin{proof}
Put $\mathbf w=z_2-z_1$ and let $J_1$ be the matrix $J$ of \eqref{eq:adjointmatrices} along $z_1$. Then $\mathbf w'=\mathbf f(z_2)-\mathbf f(z_1)=J_1\mathbf w+(0,0,x_1\Rem(U_2-U_1))^{\mathsf T}$, because the only nonlinear term of $\mathbf f$ is $ke^{U}$ in the third component and $ke^{U_2}-ke^{U_1}=x_1(U_2-U_1)+x_1\Rem(U_2-U_1)$. By \eqref{att:eq:normal}, and since the third component of $n$ is $-1$,
$\Delta'=n'\cdot\mathbf w+n\cdot J_1\mathbf w-x_1\Rem(U_2-U_1)=-(1+\eta)\Delta-x_1\Rem(U_2-U_1)$. For $s<t$,
\[
\Delta(t)=\Delta(s)\,e^{-\int_s^t(1+\eta)\,\mathrm dr}-\int_s^te^{-\int_r^t(1+\eta)\,\mathrm dr'}\,x_1(r)\,\Rem\bigl(U_2(r)-U_1(r)\bigr)\,\mathrm dr .
\]
$\Delta$ is bounded and $1+\eta\ge1$; letting $s\to-\infty$ gives
\[
\Delta(t)=-\int_{-\infty}^te^{-\int_r^t(1+\eta)\,\mathrm dr'}\,x_1(r)\,\Rem\bigl(U_2(r)-U_1(r)\bigr)\,\mathrm dr\le0 .
\]
If $\Delta(t_0)=0$, the integrand vanishes on $(-\infty,t_0]$; as $x_1>0$ and $\Rem(d)>0$ for $d\ne0$, $U_2\equiv U_1$ there, hence $p_2=U_2'=p_1$ and $y_2=k-\mu p_2-p_2'=y_1$ there, and $z_2\equiv z_1$ by uniqueness, a contradiction.
\end{proof}
Theorem~\ref{thm:support} is the case in which $z_1$ is a cycle or the equilibrium.
\begin{proof}[Proof of (A3)]
For $\mathbf a\ne\mathbf b$ in $\cA$ apply Theorem~\ref{thm:twoorbit} to $z_1=\varphi_{\cdot}(\mathbf a)$, $z_2=\varphi_{\cdot}(\mathbf b)$, which are distinct as functions also when $\mathbf b$ lies on the orbit of $\mathbf a$, and evaluate at $t=0$.
\end{proof}

\begin{lemma}[Gauges on the attractor]\label{att:lem:gauges}
On $\cA$:
\begin{enumerate}[label=\textup{(\roman*)}]
\item $O:=n\cdot\mathbf f=\zeta p-\eta(k-\mu p-y)-(ke^{U}-y)=0$;
\item $K=\zeta U-\eta p-y+k>0$ off $E$, and $K'+(1+\eta)K=k\,\psiU(U)$;
\item $W=\zeta'p-\eta'p'>0$ off $E$, and $W'+(1+\eta)W=xp^2$;
\item $\LamE=(y-k)-\zeta_EU+\eta_Ep>0$ off $E$.
\end{enumerate}
\end{lemma}
\begin{proof}
(i) Along an orbit $z(\cdot)$ in $\cA$, $z''=Jz'$, so by \eqref{att:eq:normal} $O'=(n'+J^{\mathsf T}n)\cdot z'=-(1+\eta)O$. $O$ is bounded on $\R$ and $O(t)=O(s)e^{-\int_s^t(1+\eta)\,\mathrm dr}\to0$ as $s\to-\infty$.

(ii) $K(\mathbf a)=-n(\mathbf a)\cdot(E-\mathbf a)>0$ by \eqref{att:eq:support} with $\mathbf b=E$, and the equation is \eqref{att:eq:deltaode} with $z_2\equiv E$, for which $\Delta=-K$ and $x\,\Rem(-U)=ke^{U}(e^{-U}-1+U)=k\,\psiU(U)$.

(iv) This is \eqref{att:eq:support} with $\mathbf a=E$.

(iii) By \eqref{att:eq:ric}, $\zeta''=xp-\eta'\zeta-(1+\eta)\zeta'$ and $\eta''=\zeta'+(\mu-1)\eta'-2\eta\eta'$. Also $p''=-\mu p'-(x-y)$, and (i) in the form $x-y=\zeta p-\eta p'$ gives $p''=(\eta-\mu)p'-\zeta p$. Hence
\begin{align*}
W'&=\zeta''p+\zeta'p'-\eta''p'-\eta'p''\\
&=\bigl(xp-\eta'\zeta-(1+\eta)\zeta'\bigr)p-\bigl((\mu-1)\eta'-2\eta\eta'\bigr)p'-\eta'\bigl((\eta-\mu)p'-\zeta p\bigr)\\
&=xp^2-(1+\eta)\zeta'p+(1+\eta)\eta'p'=xp^2-(1+\eta)W ,
\end{align*}
the terms $\zeta'p'$ having cancelled in the second line. Without the constraint $O=0$ the residual of this identity is $-\eta'O$. $W$ is bounded on $\R$, so $W(t)=\int_{-\infty}^te^{-\int_r^t(1+\eta)\,\mathrm dr'}x(r)p(r)^2\,\mathrm dr\ge0$, with equality only if $p\equiv0$ on $(-\infty,t]$, which forces $y\equiv k$ and $x\equiv k$: the orbit is $E$.
\end{proof}

\subsection{The graph and the reduced flow: proof of (A4) and (A5)}\label{att:sec:graph}
\begin{proof}[Proof of (A4)]
By (A2), $\mathbf a\mapsto(y_{\mathbf a},\zeta(\mathbf a),\eta(\mathbf a),\pi(\mathbf a))$ is continuous on the compact set $\cA$; so $H$ is a maximum of affine functions with slopes in a compact set: it is convex, finite, and Lipschitz with constant $\max_\cA|(\zeta,\eta)|$. For $\mathbf b\in\cA$ the term with $\mathbf a=\mathbf b$ equals $y_{\mathbf b}$ at $\pi(\mathbf b)$, and every other term is smaller than $y_{\mathbf b}$ there by \eqref{att:eq:support}; so $H(\pi(\mathbf b))=y_{\mathbf b}$. Hence $\pi|_\cA$ is injective, $\cA$ is the graph of $h=H|_{\Dsc}$, and $\pi|_\cA$, a continuous bijection of a compact space onto $\Dsc$, is a homeomorphism; $\Dsc$ is compact and connected by (A1). So $g=(\zeta,-\eta)\circ(\pi|_\cA)^{-1}$ is continuous on $\Dsc$. For $q=\pi(\mathbf a)$, $q'=\pi(\mathbf b)$, \eqref{att:eq:support} reads $h(q')-h(q)\ge g(q)\cdot(q'-q)$; exchanging $\mathbf a$ and $\mathbf b$, $h(q)-h(q')\ge g(q')\cdot(q-q')$. Together these are \eqref{att:eq:whitney}. Therefore, with $\omega_g$ the modulus of continuity of $g$ on $\Dsc$,
\begin{equation}\label{att:eq:jet}
\bigl|h(q')-h(q)-g(q)\cdot(q'-q)\bigr|\le\omega_g(|q'-q|)\,|q'-q|\qquad(q,q'\in\Dsc).
\end{equation}
This is the hypothesis of (I5) for the jet $(h,g)$ on $\Dsc$, and at an interior point $q$ it is differentiability of $h$ with $\nabla h(q)=g(q)$. Continuity of $g$ gives $h\in C^1(\inter\Dsc)$; the sign statements are \eqref{att:eq:bounds}.
\end{proof}
\begin{proof}[Proof of (A5)]
$F$ is globally Lipschitz because $H$ is, so $\chi_t$ is a complete flow on $\R^2$. For $q\in\Dsc$ let $\mathbf a=(q,h(q))\in\cA$ and $z(t)=\varphi_t(\mathbf a)\in\cA$. Then $y(t)=H(U(t),p(t))$, so $(U,p)$ solves $q'=F(q)$; by uniqueness $\pi(\varphi_t\mathbf a)=\chi_t(q)$ for all $t\in\R$, and $\Dsc$ is invariant. Equation \eqref{att:eq:pde} is Lemma~\ref{att:lem:gauges}(i) written at the point $(u,v,h(u,v))$; on $\inter\Dsc$, $(\zeta,-\eta)=\nabla h$. On $\inter\Dsc$, $F=(v,k-\mu v-h)$ with $h\in C^1$, and $\diver F=\partial_v(k-\mu v-h)=-\mu+\eta$. The statements on $K$ and $W$ are Lemma~\ref{att:lem:gauges}(ii),(iii).
\end{proof}

\begin{lemma}[The reduced flow on the $u$-axis]\label{att:lem:axis}
For $(u,0)\in\Dsc$,
\begin{equation}\label{att:eq:axis}
h(u,0)-k=\frac{k(e^{u}-1)}{1+\eta(u,0)},\qquad\text{so that}\qquad v'=-\frac{k(e^{u}-1)}{1+\eta(u,0)}\quad\text{at }(u,0).
\end{equation}
The reduced flow crosses the positive $u$-axis downwards and the negative $u$-axis upwards, and $(0,0)$ is the only equilibrium of $F$ in $\Dsc$.
\end{lemma}
\begin{proof}
At $v=0$, \eqref{att:eq:pde} reads $-\eta(k-h)=ke^{u}-h$, that is $(1+\eta)(h-k)=k(e^{u}-1)$, and $v'=k-h$ there. If $F(q)=0$ with $q\in\Dsc$, then $\chi_tq=q$, so $(q,h(q))$ is an equilibrium of \eqref{eq:logsystem}, that is $E$.
\end{proof}
On cycles, \eqref{att:eq:axis} is the first identity of \eqref{eq:crestidentities}.

\begin{remark}[The gauges in terms of $h$]\label{att:rem:convexgauges}
On $\inter\Dsc$, $K=k-[h-uh_u-vh_v]$ is the height of $E$ above the tangent plane of the graph at $(u,v)$, and $W=\frac{\mathrm d}{\mathrm dt}\bigl(\nabla h(\chi_tq)\bigr)\cdot F$ measures the change of the supporting gradient in the velocity direction. Where $h$ is twice differentiable, $W=F^{\mathsf T}(\nabla^2h)\,F$ is its Hessian quadratic form, not the complete second time derivative along a curved trajectory:
\[
\frac{\mathrm d^2}{\mathrm dt^2}h(\chi_tq)
=F^{\mathsf T}(\nabla^2h)F+\nabla h\cdot DF\,F.
\]
The quantities $K$ and $W$ express convexity of the graph.
\end{remark}

\begin{proposition}[Radial scaling]\label{att:prop:scaling}
For $q=(u,v)\in\R^2$ and $\theta\ge0$ put $\mathsf b_\theta(q)=\theta\bigl(H(q)-k\bigr)+k-H(\theta q)$, and write $(a_1,a_2)\wedge(b_1,b_2)=a_1b_2-a_2b_1$. Then
\begin{equation}\label{att:eq:scaling}
F(\theta q)\wedge F(q)=-v\,\mathsf b_\theta(q).
\end{equation}
For $q\in\Dsc$ write $K(q)$ for the value of $K$ at the point $(q,h(q))$ of $\cA$.
\begin{enumerate}[label=\textup{(\roman*)}]
\item If $0\le\theta\le1$, then $\mathsf b_\theta(q)\ge0$ for all $q\in\R^2$. If moreover $0<\theta<1$ and $\theta q\in\Dsc\setminus\{0\}$, then $\mathsf b_\theta(q)\ge(1-\theta)K(\theta q)>0$.
\item If $q\in\Dsc$ and $\theta\ge0$, then $\mathsf b_\theta(q)\le(1-\theta)K(q)$. In particular $\mathsf b_\theta(q)<0$ for $\theta>1$ and $q\in\Dsc\setminus\{0\}$.
\end{enumerate}
\end{proposition}
\begin{proof}
The identity follows by expanding $F(\theta q)=(\theta v,k-\mu\theta v-H(\theta q))$ and $F(q)=(v,k-\mu v-H(q))$; the terms $\mu\theta v^2$ cancel. Since $H(0,0)=y_E=k$ by (A4), $\mathsf b_\theta(q)=\theta H(q)+(1-\theta)H(0)-H(\theta q+(1-\theta)0)$, which is nonnegative for $0\le\theta\le1$ by convexity of $H$.

By definition of $H$ and by (A4), for every $q_1\in\Dsc$ and every $q'\in\R^2$,
\begin{equation}\label{att:eq:plane}
H(q')\ge h(q_1)+g(q_1)\cdot(q'-q_1),\qquad K(q_1)=g(q_1)\cdot q_1-h(q_1)+k .
\end{equation}
(i) Let $0<\theta<1$ and $q_1=\theta q\in\Dsc\setminus\{0\}$. Taking $q'=q$ in \eqref{att:eq:plane}, $H(q)\ge h(q_1)+(1-\theta)g(q_1)\cdot q$ and $\theta\,g(q_1)\cdot q=K(q_1)-k+h(q_1)$, so
\[
\mathsf b_\theta(q)\ge\theta h(q_1)+(1-\theta)\bigl(K(q_1)-k+h(q_1)\bigr)+(1-\theta)k-h(q_1)=(1-\theta)K(q_1),
\]
and $K(q_1)>0$ by (A5).
(ii) Taking $q_1=q$ and $q'=\theta q$ in \eqref{att:eq:plane}, $H(\theta q)\ge h(q)-(1-\theta)g(q)\cdot q=\theta h(q)+(1-\theta)k-(1-\theta)K(q)$, which is the claim.
\end{proof}
The inequality $\mathsf b_\theta\ge0$ therefore holds for $0\le\theta\le1$ only; for $\theta>1$ the sign is reversed on $\Dsc\setminus\{0\}$. Geometrically, let $\mathcal J$ be a projected cycle, which is traversed clockwise (Lemma~\ref{att:lem:cyclegeom}), and $\theta\mathcal J$ a homothetic copy. At the point $\theta q$ of $\theta\mathcal J$ the tangent of $\theta\mathcal J$ is parallel to $F(q)$, and the field $F(\theta q)$ points to the outer side of $\theta\mathcal J$ exactly when $F(q)\wedge F(\theta q)=v\,\mathsf b_\theta(q)>0$. Where $v>0$ the field therefore crosses the shrunk copies $(\theta<1$, $\theta q\in\Dsc)$ outwards and the enlarged copies inwards, towards $\theta=1$ in both cases; where $v<0$ it crosses them away from $\theta=1$. For $\theta q\notin\Dsc$ the field $F$ is the Lipschitz extension and does not represent the flow of \eqref{eq:logsystem}. Convexity of $H$ thus yields a comparison with scaled orbits whose direction reverses with the sign of $v$; it does not yield a one-signed comparison on an annulus.

\subsection{Above the Hopf value: proof of (A6) and (A7)}\label{att:sec:above}
Assume $k>\kH$. The characteristic polynomial of $J_E$, the matrix $J$ at $E$, is $z^3+Az^2+\mu z+k$. It has the real root $-\gamma_E<-A$: the cubic $\gamma(\gamma-1)(\gamma-\mu)$ equals $\kH<k$ at $\gamma=A$ and increases for $\gamma>\max(1,\mu)$. The remaining two roots have sum $\gamma_E-A>0$ and product $k/\gamma_E>0$, and they cannot be nonnegative reals because the coefficients are positive; so they form a complex pair with positive real part. Thus $E$ is hyperbolic, with a one-dimensional stable manifold tangent to the eigenvector $\mathbf e_E=(1,-\gamma_E,-\gamma_E(\gamma_E-\mu))^{\mathsf T}\in\inter C_+$ and a two-dimensional unstable manifold tangent to the real plane $E^u$. By \eqref{att:eq:normal} with the constant vector $n_E=(\zeta_E,-\eta_E,-1)^{\mathsf T}$, $J_E^{\mathsf T}n_E=-\gamma_En_E$; a left eigenvector for the simple eigenvalue $-\gamma_E$ annihilates the other eigenvectors, so
\begin{equation}\label{att:eq:Eu}
E^u=n_E^{\perp}=\{\mathbf v=(v_1,v_2,v_3)\in\R^3:\ v_3=\zeta_E\,v_1-\eta_E\,v_2\},
\end{equation}
a graph over the $(U,p)$-plane.

\begin{proof}[Proof of (A6)]
\emph{Cone exclusion.} Let $\mathbf a\in\cA\setminus\{E\}$ and $\mathbf a-E=(U,p,y-k)$. If $\mathbf a-E\in C_+$ then $\zeta_EU-\eta_Ep-(y-k)\ge0$, contradicting Lemma~\ref{att:lem:gauges}(iv). If $\mathbf a-E\in C_-$ then $K=\zeta U-\eta p-(y-k)\le0$, because $\zeta,\eta>0$, contradicting Lemma~\ref{att:lem:gauges}(ii).

\emph{No forward convergence to $E$.} If $\mathbf a\in\cA\setminus\{E\}$ and $\varphi_t(\mathbf a)\to E$, then by (I3a) $\varphi_t(\mathbf a)$ lies on the local stable manifold for large $t$, a $C^1$ curve tangent to $\mathbf e_E$. Hence the limit points of $(\varphi_t(\mathbf a)-E)/|\varphi_t(\mathbf a)-E|$ lie in $\{\pm\mathbf e_E/|\mathbf e_E|\}\subset\inter C_+\cup\inter C_-$, and $\varphi_t(\mathbf a)\in E+C_\pm$ for large $t$, contradicting cone exclusion.

\emph{$W^u(E)\subset\cA$ and $(0,0)\in\inter\Dsc$.} A point of $W^u(E)$ has a backward orbit converging to $E$, hence bounded and, near $E$, inside the open set $\Dom$, and a forward orbit which stays in $\Dom$ and is bounded by Lemma~\ref{lem:absorbing}; by (A1) it lies in $\cA$. The local unstable manifold is a $C^1$ surface tangent at $E$ to $E^u$, on which $\mathrm d\pi$ is injective by \eqref{att:eq:Eu}; by the inverse function theorem $\pi(W^u_{\rm loc}(E))$ contains a neighbourhood $N_E$ of $(0,0)$ in $\R^2$. So $N_E\subset\Dsc$, and every point of $N_E\setminus\{0\}$ has $\alpha$-limit set $\{0\}$ under $\chi_t$; in particular $N_E$ contains no periodic point.

\emph{Limit sets.} Let $q\in\Dsc\setminus\{0\}$. Its $\chi_t$-orbit lies in the compact invariant set $\Dsc$, in which $0$ is the only equilibrium of $F$ (Lemma~\ref{att:lem:axis}). By (I2a), $\omega(q)$ is $\{0\}$, or a periodic orbit, or contains an orbit other than $\{0\}$ in $\Dsc$ converging to $0$ as $t\to+\infty$. The first and third alternatives contradict the second step, after lifting to $\cA$ by (A5). So $\omega(q)$ is a periodic orbit of $\chi_t$ in $\Dsc$, that is, the projection of a cycle. Applying (I2a) to $-F$: $\alpha(q)$ is $\{0\}$, or a periodic orbit, or contains an orbit other than $\{0\}$ in $\Dsc$ which converges to $0$ as $t\to\pm\infty$; the last alternative is excluded in the same way.

\emph{$\alpha(q)\neq\omega(q)$ for non-periodic $q$.} Suppose $\alpha(q)=\omega(q)=\mathcal J$ is a periodic orbit and $q\notin\mathcal J$. Fix $q_\ast\in\mathcal J$. Since $F(q_\ast)\ne0$ and $F$ is continuous, a sufficiently short closed segment $I$ with midpoint $q_\ast$, perpendicular to $F(q_\ast)$, is a transversal in the sense of (I2b): $F$ does not vanish on $I$ and is nowhere parallel to $I$. The component of $F$ normal to $I$ does not vanish on $I$, so it has a fixed sign on a neighbourhood of $I$, and the signed distance to the line containing $I$ is strictly monotone along orbit arcs in that neighbourhood. Hence every crossing of $I$ is isolated in time, and every orbit that passes close enough to $q_\ast$ crosses $I$ close to $q_\ast$ within a short time. As $q_\ast\in\omega(q)\cap\alpha(q)$, the orbit of $q$ meets $I$ at times tending to $+\infty$ and at times tending to $-\infty$, in points converging to $q_\ast$. Its crossing times form a closed discrete set which is unbounded in both directions, that is, a bi-infinite increasing sequence, and by (I2b) the corresponding points form a monotone sequence along $I$; it is strictly monotone because $q$ is not periodic. A monotone sequence in $I$ has a limit at each of its two ends, and here both limits equal $q_\ast$, since a subsequence converges to $q_\ast$ at each end. A strictly monotone bi-infinite sequence cannot have the same limit at both ends. This contradiction proves the claim.

\emph{Existence.} For $q\in N_E\setminus\{0\}$, $\omega(q)$ is the projection of a cycle.
\end{proof}

\begin{lemma}[Geometry of projected cycles]\label{att:lem:cyclegeom}
Let $\mathcal J\subset\Dsc$ be a periodic orbit of $\chi_t$; equivalently, by (A5), the projection of a cycle. Then $\mathcal J$ is a Jordan curve with $0$ in its interior, traversed clockwise, and it meets the positive $u$-axis in exactly one point $(\xi,0)$, its \emph{crest}, and the negative $u$-axis in exactly one point. The points $(u,0)$ with $0<u<\xi$ lie in $\operatorname{ins}\mathcal J$ and those with $u>\xi$ in $\operatorname{ext}\mathcal J$. Two such orbits $\mathcal J,\mathcal J'$ with crests $\xi<\xi'$ are disjoint and satisfy $\mathcal J\subset\operatorname{ins}\mathcal J'$, $\mathcal J'\subset\operatorname{ext}\mathcal J$.
\end{lemma}
\begin{proof}
$\mathcal J$ is a Jordan curve not through the equilibrium $0$. Since $u'=v$, the periodic function $u$ has extrema, at which $v=0$. By Lemma~\ref{att:lem:axis}, every crossing of $\{v=0,u>0\}$ is transverse and downward, every crossing of $\{v=0,u<0\}$ is transverse and upward, and $\mathcal J$ does not meet $\{v=0,u=0\}$. At the maximum of $u$ along $\mathcal J$, $v=0$ and $v'\le0$, so $u_{\max}>0$; likewise $u_{\min}<0$. Parametrise $\mathcal J$ by the flow over one period; this is a homeomorphism of the circle onto $\mathcal J$. The crossings of the ray $\{v=0,u>0\}$ are transverse, hence isolated in time, so there are finitely many in one period, say $N$ of them, and $N\ge1$ because $\mathcal J$ crosses the ray at the maximum of $u$. Between two consecutive crossings the orbit stays in the plane slit along this ray, where the polar angle about $0$ has a continuous branch with values in $(0,2\pi)$; since all crossings are downward, this branch is close to $2\pi$ just after a crossing and close to $0$ just before the next one. A continuous determination of the angle along $\mathcal J$ therefore decreases by $2\pi$ from each crossing to the next, and the winding number of $\mathcal J$ about $0$ is $-N$. By (I4) the winding number of a Jordan curve about a point off the curve is $0$ or $\pm1$. So $N=1$, the winding number is $-1$ (clockwise), and $0\in\operatorname{ins}\mathcal J$. The same argument, with the branch of the angle in $(-\pi,\pi)$ and upward crossings, shows that $\mathcal J$ crosses the negative ray exactly once. The open segment from $0$ to $(\xi,0)$ does not meet $\mathcal J$ and contains points near $0$, so it lies in $\operatorname{ins}\mathcal J$; the ray beyond $(\xi,0)$ does not meet $\mathcal J$ and is unbounded, so it lies in $\operatorname{ext}\mathcal J$. Distinct periodic orbits are disjoint; if $\xi<\xi'$, the point $(\xi,0)\in\mathcal J$ is inside $\mathcal J'$ and the point $(\xi',0)\in\mathcal J'$ is outside $\mathcal J$, so by connectedness $\mathcal J\subset\operatorname{ins}\mathcal J'$ and $\mathcal J'\subset\operatorname{ext}\mathcal J$.
\end{proof}
Lemma~\ref{att:lem:cyclegeom} agrees with Theorem~\ref{thm:geometry}; it is proved here from the reduced flow so that the argument also applies to the comparison of a cycle with a non-periodic orbit of $\Dsc$. Write $\mathcal J_\xi$ for the projected cycle with crest $\xi$, and $\mathrm{Cr}\subset(0,\infty)$ for the set of crests.

\begin{lemma}\label{att:lem:crests}
$\mathrm{Cr}$ is nonempty and compact in $(0,\infty)$.
\end{lemma}
\begin{proof}
It is nonempty by (A6), bounded since $\Dsc$ is, and bounded away from $0$ since $N_E$ contains no periodic point. Let $\xi_j\in\mathrm{Cr}$, $\xi_j\to\bar\xi$, and $\bar q=(\bar\xi,0)\in\Dsc\setminus\{0\}$. Suppose $\bar q$ is not periodic; then $\xi_j\ne\bar\xi$ for all $j$. By (A6), $\omega(\bar q)=\mathcal J_{\xi_\omega}$, and $\alpha(\bar q)$ is $\{0\}$, in which case we put $\xi_\alpha=0$ and $\mathcal J_0:=\{0\}$, or $\alpha(\bar q)=\mathcal J_{\xi_\alpha}$; in both cases $\xi_\alpha\ne\xi_\omega$. Let $\xi_1<\xi_2$ be these two numbers. The orbit of $\bar q$ is disjoint from $\mathcal J_{\xi_1}\cup\mathcal J_{\xi_2}$ and accumulates on both, so it lies in the open annulus between them: inside $\mathcal J_{\xi_1}$ it could not accumulate on $\mathcal J_{\xi_2}\subset\operatorname{ext}\mathcal J_{\xi_1}$, and outside $\mathcal J_{\xi_2}$ it could not accumulate on $\mathcal J_{\xi_1}$. Hence $\xi_1<\bar\xi<\xi_2$ by Lemma~\ref{att:lem:cyclegeom}, and for large $j$ also $\xi_1<\xi_j<\xi_2$, so $\mathcal J_{\xi_j}$ separates $\mathcal J_{\xi_1}$ from $\mathcal J_{\xi_2}$. The orbit of $\bar q$ is disjoint from the periodic orbit $\mathcal J_{\xi_j}$, hence lies inside it or outside it, and then cannot accumulate on $\mathcal J_{\xi_2}$, respectively on $\mathcal J_{\xi_1}$. This contradiction shows that $\bar q$ is periodic, so $\bar\xi\in\mathrm{Cr}$.
\end{proof}
Let $\xi_{\rm out}=\max\mathrm{Cr}$ and $\xi_{\rm in}=\min\mathrm{Cr}$, with cycles $\Gamma_{\rm out},\Gamma_{\rm in}$ and projections $\mathcal J_{\rm out}=\pi(\Gamma_{\rm out})$, $\mathcal J_{\rm in}=\pi(\Gamma_{\rm in})$. Let $\Omega=\operatorname{ins}\mathcal J_{\rm out}$ and $\overline\Omega=\Omega\cup\mathcal J_{\rm out}$. By Lemma~\ref{att:lem:cyclegeom} every projected cycle lies in $\overline\Omega$.

\begin{proof}[Proof of (A7)]
\emph{$\Dsc\subset\overline\Omega$.} Let $q\in\Dsc\setminus\overline\Omega$. Its orbit stays in $\operatorname{ext}\mathcal J_{\rm out}$, so $\alpha(q)$ and $\omega(q)$ lie in $\Dsc\setminus\Omega$. By (A6) both are projected cycles, since $\{0\}$ is excluded by $0\in\Omega$; every projected cycle lies in $\overline\Omega$, so both equal $\mathcal J_{\rm out}$, and $q$ is not periodic. This contradicts (A6).

\emph{$\overline\Omega\subset\Dsc$.} Suppose $q_0\in\Omega\setminus\Dsc$. Let $y_{\min}=\min_\cA y>0$ and $H_0=\max(H,y_{\min})$ on $\overline\Omega$; then $H_0=h$ on $\mathcal J_{\rm out}$, and $\mathbf S_0=\{(q,H_0(q)):q\in\overline\Omega\}$ is a compact subset of $\Dom$: its points have $y>0$, and $p<\lambda$ on $\overline\Omega$ because the linear function $v$ attains its maximum over $\overline\Omega$ on $\mathcal J_{\rm out}\subset\Dsc$. For $t\ge0$ define $\Theta_t:\overline\Omega\to\R^2$ by $\Theta_t(q)=\pi\varphi_t(q,H_0(q))$; the map $(t,q)\mapsto\Theta_t(q)$ is continuous on $[0,\infty)\times\overline\Omega$. On $\mathcal J_{\rm out}$, $\Theta_t=\chi_t$ maps $\mathcal J_{\rm out}$ onto itself, and $\Theta_0$ is the identity. Fix $t$ and suppose $q_0\notin\Theta_t(\overline\Omega)$. By (I4) there is a homeomorphism $\Upsilon$ of the closed unit disc onto $\overline\Omega$ mapping the unit circle $S^1$ onto $\mathcal J_{\rm out}$. The loops $\Theta_s\circ\Upsilon|_{S^1}$, $0\le s\le t$, lie in $\mathcal J_{\rm out}\subset\R^2\setminus\{q_0\}$ and form a homotopy from $\Upsilon|_{S^1}$, which has winding number $\pm1$ about $q_0$, to $\Theta_t\circ\Upsilon|_{S^1}$. The latter loop extends to the continuous map $\Theta_t\circ\Upsilon$ of the closed disc into $\R^2\setminus\{q_0\}$ and so, by the last statement of (I4), has winding number $0$ about $q_0$. This contradicts the homotopy invariance of the winding number. Hence for every $j\in\mathbb N$ there is $\mathbf z_j\in\varphi_j(\mathbf S_0)$ with $\pi(\mathbf z_j)=q_0$. By (A1), $\operatorname{dist}(\mathbf z_j,\cA)\to0$; a limit point $\mathbf z\in\cA$ has $\pi(\mathbf z)=q_0$, so $q_0\in\Dsc$, a contradiction.

Thus $\Dsc=\overline\Omega$, $\partial\Dsc=\mathcal J_{\rm out}$, and $\cA=(\pi|_\cA)^{-1}(\overline\Omega)$ is a closed topological disc with boundary $\Gamma_{\rm out}$. The statement on crests is Lemma~\ref{att:lem:crests}.
\end{proof}

\subsection{Multipliers and the parameter-wise criterion: proof of (A8)}\label{att:sec:mult}
Let $\Gamma$ be a cycle, $\mathbf a_0\in\Gamma$, $T$ its period and $\mathbf M$ the monodromy matrix at $\mathbf a_0$. By \eqref{att:eq:normal}, $\tilde n(t)=e^{\int_0^t(1+\eta)\,\mathrm dr}n(t)$ solves the adjoint equation $\tilde n'=-J^{\mathsf T}\tilde n$ along $\Gamma$, so $\tilde n\cdot\mathbf v$ is constant for every solution $\mathbf v(t)$ of the variational equation $\mathbf v'=J\mathbf v$; since $n$ is periodic this gives $\mathbf M^{\mathsf T}n(\mathbf a_0)=m_s\,n(\mathbf a_0)$ with $m_s:=e^{-(1+\avg\eta)T}<1$. So $m_s$ is an eigenvalue of $\mathbf M$; another is $1$, with eigenvector $\mathbf f(\mathbf a_0)$, and by Liouville's formula the product of the three is $e^{-AT}$, since $\tr J=-A$. Hence the third is $m_t=e^{-AT}/m_s=e^{(\avg\eta-\mu)T}$; this is Proposition~\ref{prop:multipliers}. Since $\eta>\max(0,\mu-1)$ by \eqref{att:eq:bounds}, $2\avg\eta>\mu-1$, that is $m_s<m_t$; also $m_s<1$. So $m_s$ is simple and
\begin{equation}\label{att:eq:nperp}
n(\mathbf a_0)^\perp=\text{generalised eigenspace of }\mathbf M\text{ for }\{1,m_t\}\ni\mathbf f(\mathbf a_0),\qquad n(\mathbf a_0)\cdot e_s\ne0,
\end{equation}
where $e_s$ is an eigenvector of $\mathbf M$ for $m_s$.

\begin{proof}[Proof of (A8)]
(a) By (A5), $\pi(\Gamma)$ is a periodic orbit of $\chi_t$ with period $T$. If it lies in $\inter\Dsc$, $F$ is $C^1$ near it; the planar variational equation has the periodic solution $F(\chi_tq)$, so one of its multipliers is $1$ and, by Liouville's formula, the other is $\exp\int_0^T\diver F\,\mathrm dt=\exp\int_0^T(\eta-\mu)\,\mathrm dt=m_t$, $\eta$ on $\Gamma$ being the function of (A2). For $\Gamma_{\rm out}$: by \eqref{att:eq:jet} and (I5) there is $\widehat h\in C^1(\R^2)$ with $\widehat h=h$ and $\nabla\widehat h=g$ on $\Dsc$. Then $\widehat F=F$ on $\Dsc$, each $\chi_t$-orbit in $\Dsc$ is an orbit of $\widehat F$ by uniqueness for the $C^1$ field $\widehat F$, and $\diver\widehat F=\eta-\mu$ on $\Dsc$; the same computation therefore gives the multiplier $m_t$ for $\mathcal J_{\rm out}$ as a periodic orbit of $\widehat F$, whatever the extension. By (A7) every other projected cycle lies in $\Omega=\inter\Dsc$.

(b) Let $m_t>1$ and suppose $\mathbf b\in\cA\setminus\Gamma$ has $\omega(\mathbf b)=\Gamma$. Take a small disc $\mathcal T$ transverse to $\mathbf f$ at $\mathbf a_0$, with first-return map $\Pi$; its derivative $\mathrm d\Pi(\mathbf a_0)$ has the eigenvalues $m_s<1<m_t$, with stable eigenvector $\bar e_s=e_s+r\,\mathbf f(\mathbf a_0)$ for some $r\in\R$. The orbit of $\mathbf b$ meets $\mathcal T$ in points $\mathbf b_j=\Pi^j(\mathbf b_0)\to\mathbf a_0$, $\mathbf b_j\ne\mathbf a_0$, which by (I3b) lie, for all large $j$, on the local stable manifold of $\mathbf a_0$ for $\Pi$, a $C^1$ curve tangent to $\bar e_s$. Hence the limit points of $(\mathbf b_j-\mathbf a_0)/|\mathbf b_j-\mathbf a_0|$ lie in $\{\pm\bar e_s/|\bar e_s|\}$. But $\mathbf b_j,\mathbf a_0\in\cA$, so by \eqref{att:eq:whitney} and \eqref{att:eq:jet}
\[
\bigl|n(\mathbf a_0)\cdot(\mathbf b_j-\mathbf a_0)\bigr|=h(\pi\mathbf b_j)-h(\pi\mathbf a_0)-g(\pi\mathbf a_0)\cdot(\pi\mathbf b_j-\pi\mathbf a_0)\le\omega_g(|\mathbf b_j-\mathbf a_0|)\,|\mathbf b_j-\mathbf a_0| ,
\]
so $n(\mathbf a_0)\cdot\bar e_s=0$. By \eqref{att:eq:nperp} and Lemma~\ref{att:lem:gauges}(i), $n(\mathbf a_0)\cdot\bar e_s=n(\mathbf a_0)\cdot e_s\ne0$, a contradiction. If $m_t<1$, then $\Gamma$ is orbitally stable by (I3c): every neighbourhood $V$ of $\Gamma$ contains a neighbourhood $V'$ whose forward orbits stay in $V$. If $\alpha(\mathbf b)=\Gamma$ then $\varphi_{-t_j}(\mathbf b)\in V'$ for some $t_j\to\infty$, so $\mathbf b\in V$ for every $V$, that is $\mathbf b\in\Gamma$.

(c) Let $q\in N_E\setminus\{0\}$, so that $\alpha(q)=\{0\}$ by (A6). The orbit of $q$ is not periodic, hence disjoint from $\mathcal J_{\rm in}$, and it has the point $0\in\operatorname{ins}\mathcal J_{\rm in}$ in its closure; so it lies in $\operatorname{ins}\mathcal J_{\rm in}$. By (A6), $\omega(q)$ is a projected cycle contained in the closure of $\operatorname{ins}\mathcal J_{\rm in}$; its crest is at most $\xi_{\rm in}$, so $\omega(q)=\mathcal J_{\rm in}$. Thus an orbit in $\cA\setminus\Gamma_{\rm in}$ has $\omega$-limit set $\Gamma_{\rm in}$, and (b) gives $m_t(\Gamma_{\rm in})\le1$. Suppose $m_t(\Gamma_{\rm out})>1$. Then $\Gamma_{\rm out}$ is hyperbolic with a two-dimensional local unstable manifold $W^u_{\rm loc}$ as in (I3b): a $C^1$ surface containing $\Gamma_{\rm out}$ whose tangent plane at $\mathbf a_0\in\Gamma_{\rm out}$ is spanned by $\mathbf f(\mathbf a_0)$ and the eigenvector for $m_t$, and so equals $n(\mathbf a_0)^\perp$ by \eqref{att:eq:nperp}. Since $n$ has third component $-1$, $\mathrm d\pi$ is injective on this plane, and by the inverse function theorem $\pi(W^u_{\rm loc})$ contains a neighbourhood of $\mathcal J_{\rm out}$ in $\R^2$. After shrinking, $W^u_{\rm loc}$ lies in the open set $\Dom$, which contains $\Gamma_{\rm out}$. Its points have backward orbits converging to $\Gamma_{\rm out}$, hence bounded, and forward orbits which are bounded by Lemma~\ref{lem:absorbing}; so $W^u_{\rm loc}\subset\cA$ by (A1), and $\Dsc$ contains points outside $\mathcal J_{\rm out}$, contradicting (A7).

(d) Suppose every cycle has $m_t<1$ and there are at least two cycles. We first show that each crest is isolated in $\mathrm{Cr}$. Let $\xi\in\mathrm{Cr}$, with cycle $\Gamma_\xi$ of period $T$ and crest point $\mathbf a_\xi=(\xi,0,h(\xi,0))$. Since $p'=k-y=-g_\ast(\xi)\ne0$ at $\mathbf a_\xi$ by Lemma~\ref{att:lem:axis}, the plane $\{p=0\}$ is transverse to $\mathbf f$ there. The implicit function theorem gives a neighbourhood $\mathcal T$ of $\mathbf a_\xi$ in $\{p=0\}$, a $C^1$ function $t_\Pi:\mathcal T\to\R$ with $t_\Pi(\mathbf a_\xi)=T$ and $\varphi_{t_\Pi(\mathbf a)}(\mathbf a)\in\{p=0\}$, and the $C^1$ return map $\Pi(\mathbf a)=\varphi_{t_\Pi(\mathbf a)}(\mathbf a)$, with $\Pi(\mathbf a_\xi)=\mathbf a_\xi$. The eigenvalues of the derivative $\mathrm d\Pi(\mathbf a_\xi)$ are $m_s$ and $m_t$, both different from $1$, so $\mathrm d\Pi(\mathbf a_\xi)-\mathrm{Id}$ is invertible and, by the inverse function theorem applied to $\mathbf a\mapsto\Pi(\mathbf a)-\mathbf a$, the point $\mathbf a_\xi$ is the only fixed point of $\Pi$ in a neighbourhood $\mathcal T'\subset\mathcal T$ of $\mathbf a_\xi$. Since $\Pi$ is continuous and $\Pi(\mathbf a_\xi)=\mathbf a_\xi$ has $U=\xi>0$, we may shrink $\mathcal T'$ so that $\Pi(\mathcal T')\subset\{p=0,\ U>0\}$. Suppose crests $\xi_j\ne\xi$ converge to $\xi$, with cycles $\Gamma_j$. Their crest points $\mathbf a_j=(\xi_j,0,h(\xi_j,0))$ converge to $\mathbf a_\xi$ by continuity of $h$, so $\mathbf a_j\in\mathcal T'$ for large $j$. The point $\Pi(\mathbf a_j)$ lies on $\Gamma_j$ and in $\{p=0,\ U>0\}$. By Lemma~\ref{att:lem:cyclegeom} the projection $\pi(\Gamma_j)$ meets the positive $u$-axis only at $(\xi_j,0)$, and $\pi$ is injective on $\cA$; hence $\Gamma_j$ meets $\{p=0,\ U>0\}$ only at $\mathbf a_j$. Therefore $\Pi(\mathbf a_j)=\mathbf a_j$: the return of $\mathbf a_j$ selected by $t_\Pi$, at a time close to $T$, is its return to its own crest, and $\mathbf a_j\ne\mathbf a_\xi$ is a second fixed point of $\Pi$ in $\mathcal T'$. This contradiction shows that $\xi$ is isolated.

A compact set of isolated points is finite, so by Lemma~\ref{att:lem:crests} there are adjacent crests $\xi_1<\xi_2$ with no crest between them. The open annulus between $\mathcal J_{\xi_1}$ and $\mathcal J_{\xi_2}$ lies in $\Dsc$ by (A7) and contains no periodic point, because by Lemma~\ref{att:lem:cyclegeom} a periodic orbit in it would have its crest in $(\xi_1,\xi_2)$. Let $q$ be a point of this annulus. Its orbit stays in the annulus, so $\alpha(q)$ lies in the closure of the annulus; by (A6) $\alpha(q)$ is $\{0\}$ or a projected cycle, and the point $0$ and all projected cycles other than $\mathcal J_{\xi_1},\mathcal J_{\xi_2}$ lie outside that closure. Hence $\alpha(q)\in\{\mathcal J_{\xi_1},\mathcal J_{\xi_2}\}$. This contradicts (b), since both cycles have $m_t<1$.

(e) If all cycles are hyperbolic, $\mathrm{Cr}=\{\xi_1<\dots<\xi_N\}$ is finite by the first part of (d), which uses only $m_t\ne1$. For $q$ in the annulus between $\mathcal J_{\xi_i}$ and $\mathcal J_{\xi_{i+1}}$, $\alpha(q)$ and $\omega(q)$ both belong to $\{\mathcal J_{\xi_i},\mathcal J_{\xi_{i+1}}\}$ as in (d), and they differ by (A6). By (b) the $\alpha$-limit cycle has $m_t>1$ and the $\omega$-limit cycle has $m_t<1$. So consecutive cycles have opposite stability, and by (c) the first and the last have $m_t<1$; hence $N$ is odd.

(f) On a cycle $\int_0^T\Dt\Lambda\,\mathrm dt=0$, so
\begin{equation}\label{att:eq:dulacmean}
\lambda_cT=\int_0^T(\Dt\Lambda+\eta-\mu)\,\mathrm dt .
\end{equation}
Under the first hypothesis this is negative for every cycle, and (d) applies. Under the second, $\lambda_c(\Gamma_{\rm in})>0$ or $\lambda_c(\Gamma_{\rm out})>0$, contradicting (c). The last statement is \eqref{att:eq:dulacmean} again.
\end{proof}

\begin{corollary}[Parameter-wise criterion]\label{cor:criterion}
Let $(\mu,k)\in\cP$.
\begin{enumerate}[label=\textup{(\roman*)}]
\item The innermost and the outermost cycle satisfy $\lambda_c(\Gamma_{\rm in})\le0$ and $\lambda_c(\Gamma_{\rm out})\le0$.
\item If every cycle at $(\mu,k)$ has $\lambda_c<0$, then there is exactly one cycle $\Gamma$. It is hyperbolic and orbitally asymptotically stable, and it is the $\omega$-limit set of every solution in $\Dom$ that does not converge to $E$, that is, of every positive orbit off the one-dimensional stable manifold of $E$.
\item If there are at least two cycles at $(\mu,k)$, then some cycle has $\lambda_c\ge0$.
\end{enumerate}
\end{corollary}
\begin{proof}
(i) and (iii) are (A8)(c),(d), and so is the uniqueness in (ii). Let $\Gamma$ be the unique cycle; its transverse multipliers satisfy $m_s<1$ and $m_t=e^{\lambda_cT}<1$, so it is hyperbolic and, by (I3c), orbitally asymptotically stable: there is an open set $V\supset\Gamma$ such that every forward orbit starting in $V$ has $\omega$-limit set $\Gamma$. Let $\mathbf z\in\Dom$ and let $L$ be its $\omega$-limit set. By Lemma~\ref{lem:absorbing} the forward orbit is eventually in the compact box $\mathcal B_0\subset\Dom$, so $L$ is nonempty, compact and invariant, and consists of log-bounded entire orbits; hence $L\subset\cA$ by (A1). If $L\cap\Gamma\ne\emptyset$, the orbit of $\mathbf z$ enters $V$, and $L=\Gamma$. Suppose $L\cap\Gamma=\emptyset$. If $L$ contained a point $\mathbf a\ne E$, then by (A6) the $\omega$-limit set of $\mathbf a$ would be a cycle, hence $\Gamma$, and it would lie in the closed invariant set $L$; this is excluded. So $L=\{E\}$ and $\varphi_t(\mathbf z)\to E$. By (I3a) and the hyperbolicity of $E$ established in Section~\ref{att:sec:above}, the set of points whose orbits converge to $E$ is the stable manifold of $E$, which is one-dimensional.
\end{proof}

\begin{remark}[What is and is not proved]\label{att:rem:status}
Theorem~\ref{thm:attractor} reduces the question of uniqueness to the planar system $u'=v$, $v'=k-\mu v-h(u,v)$ on the closed Jordan domain $\Dsc$, with $h$ the restriction of a convex function, $h_u=\zeta>0$, $h_v=-\eta<-\max(0,\mu-1)$, the first-order equation \eqref{att:eq:pde} and divergence $\eta-\mu$. By Corollary~\ref{cor:criterion}, uniqueness at $(\mu,k)$ follows from $\lambda_c<0$ on every cycle there, and conversely a unique cycle has $\lambda_c\le0$; the two statements differ only in the degenerate case $\lambda_c=0$. The sign of $\avg\eta-\mu$ on an arbitrary cycle is not decided by the results of this section.
\end{remark}

\subsection{The return map on the crest half-axis}\label{att:sec:crestaxis}
The statements of this subsection are proved, but part (iii) of Proposition~\ref{prop:crestaxis} is conditional. Its hypotheses are not established, and floating-point computations indicate that they fail at some parameters; see the paragraph after the proof.

\begin{proposition}[Crest return map]\label{prop:crestaxis}
Assume $k>\kH$. Let $g_\ast(u)=h(u,0)-k=k(e^{u}-1)/(1+\eta(u,0))$ and $G_\ast(u)=\int_0^ug_\ast(s)\,\mathrm ds$ for $0\le u\le\xi_{\rm out}$.
\begin{enumerate}[label=\textup{(\roman*)}]
\item For $0<u<\xi_{\rm out}$ the $\chi_t$-orbit of $(u,0)$ returns to the half-axis $\{v=0,u>0\}$. The first return $R_1(u)$ and the return time $T_1(u)$ are $C^1$ on $(0,\xi_{\rm out})$, $R_1$ is increasing, and its fixed points are the crests in $(0,\xi_{\rm out})$. $R_1$ and $T_1$ extend continuously to $u=\xi_{\rm out}$, with $R_1(\xi_{\rm out})=\xi_{\rm out}$ and $T_1(\xi_{\rm out})$ the period of $\Gamma_{\rm out}$.
\item With $\iota(u)=\int_0^{T_1(u)}(\eta-\mu)(\chi_t(u,0))\,\mathrm dt$ and the energy coordinate $\varepsilon=G_\ast(u)$, for $0<u<\xi_{\rm out}$,
\begin{equation}\label{att:eq:returnderivative}
R_1'(u)=\frac{g_\ast(u)}{g_\ast(R_1(u))}\,e^{\iota(u)}\qquad\text{and}\qquad
\frac{\mathrm d}{\mathrm d\varepsilon}\,G_\ast\bigl(R_1(G_\ast^{-1}(\varepsilon))\bigr)=e^{\iota(u)}.
\end{equation}
\item If $u\mapsto R_1(u)/u$ is strictly decreasing on $(0,\xi_{\rm out}]$, or if $\iota$ is strictly decreasing on $(0,\xi_{\rm out})$, equivalently if the energy return map $R_\ast=G_\ast\circ R_1\circ G_\ast^{-1}$ is strictly concave, then the cycle is unique.
\end{enumerate}
\end{proposition}
\begin{proof}
(i) By (A7) and Lemma~\ref{att:lem:cyclegeom}, $(u,0)\in\inter\Dsc\setminus\{0\}$ for $0<u<\xi_{\rm out}$, and its orbit stays in the invariant open set $\inter\Dsc$, where $F$ is $C^1$. By (A6) its $\omega$-limit set is a projected cycle $\mathcal J_\xi$, and the $u$-axis is transverse to $F$ at $(\xi,0)$ by Lemma~\ref{att:lem:axis}, so the orbit meets $\{v=0,u>0\}$ at arbitrarily large times. The first return is therefore defined; its fixed points are the periodic points on the half-axis, that is, the crests.

Let $0<u_0\le\xi_{\rm out}$, let $t_0=T_1(u_0)$, and let $(u_1,0)$ be the point of first return; for $u_0=\xi_{\rm out}$, $t_0$ is the period of $\Gamma_{\rm out}$ and $u_1=\xi_{\rm out}$. The second component of $F$ equals $-g_\ast<0$ at $(u_0,0)$ and at $(u_1,0)$ by Lemma~\ref{att:lem:axis}, and it is continuous on $\R^2$; so it is negative on neighbourhoods $N_0$ of $(u_0,0)$ and $N_1$ of $(u_1,0)$, which we take inside $\{u>0\}$, and $v$ is strictly decreasing along orbit arcs in $N_0$ or $N_1$. Choose $\vartheta>0$ so small that $\chi_t(u_0,0)\in N_0$ for $|t|\le\vartheta$ and $\chi_t(u_0,0)\in N_1$ for $|t-t_0|\le\vartheta$. The compact arc $\{\chi_t(u_0,0):\vartheta\le t\le t_0-\vartheta\}$ does not meet the closed set $\{v=0,u\ge0\}$. By continuity of the flow, for all $u$ close to $u_0$ with $u\le\xi_{\rm out}$ the arc of $(u,0)$ over $[\vartheta,t_0-\vartheta]$ does not meet this set either, $v(\chi_t(u,0))<0$ for $0<t\le\vartheta$, and $v(\chi_t(u,0))$ is strictly decreasing on $[t_0-\vartheta,t_0+\vartheta]$, positive at $t_0-\vartheta$ and negative at $t_0+\vartheta$. Hence $T_1(u)$ is the unique zero of $v(\chi_t(u,0))$ in $(t_0-\vartheta,t_0+\vartheta)$. Since $\vartheta$ is arbitrarily small, $T_1$ and $R_1$ are continuous at $u_0$. For $u_0<\xi_{\rm out}$ the implicit function theorem, applied to the $C^1$ function $(t,u)\mapsto v(\chi_t(u,0))$ at $(t_0,u_0)$, where its $t$-derivative is $-g_\ast(u_1)\ne0$, shows that $T_1$ and $R_1$ are $C^1$ near $u_0$. Monotonicity of $R_1$ follows from (ii), because $g_\ast>0$ on $(0,\xi_{\rm out}]$.

(ii) Let $q(t)=\chi_t(u,0)$ and $\jmath(t)=\partial_uq(t)$. Both $F(q(t))$ and $\jmath(t)$ solve the variational equation along $q(t)$, so the wedge $w=F(q)\wedge\jmath$ satisfies $w'=(\diver F)\,w$, and $w(0)=(0,-g_\ast(u))\wedge(1,0)=g_\ast(u)$. Differentiating $\chi_{T_1(u)}(u,0)=(R_1(u),0)$ gives $\jmath(T_1)+T_1'(u)F=(R_1'(u),0)$, and wedging with $F(q(T_1))=(0,-g_\ast(R_1(u)))$ gives $w(T_1)=g_\ast(R_1(u))R_1'(u)$. Since $\diver F=\eta-\mu$ on $\inter\Dsc$, this is the first formula. The second follows from $G_\ast'=g_\ast$.

(iii) If $R_1(u)/u$ is strictly decreasing on $(0,\xi_{\rm out}]$, then, since its value at $\xi_{\rm out}$ is $1$ by (i), $R_1(u)>u$ on $(0,\xi_{\rm out})$ and $R_1$ has no fixed point there.

Put $\varepsilon_{\max}=G_\ast(\xi_{\rm out})$. Since $g_\ast$ is continuous on $[0,\xi_{\rm out}]$ and positive on $(0,\xi_{\rm out}]$, $G_\ast$ is an increasing $C^1$ bijection of $(0,\xi_{\rm out}]$ onto $(0,\varepsilon_{\max}]$ with $C^1$ inverse. So $R_\ast=G_\ast\circ R_1\circ G_\ast^{-1}$ is positive and $C^1$ on $(0,\varepsilon_{\max})$, and by (ii) $R_\ast'(\varepsilon)=e^{\iota(G_\ast^{-1}(\varepsilon))}$. As $G_\ast^{-1}$ is increasing, $\iota$ is strictly decreasing on $(0,\xi_{\rm out})$ if and only if $R_\ast'$ is strictly decreasing on $(0,\varepsilon_{\max})$, and this holds if and only if $R_\ast$ is strictly concave: a differentiable function with strictly decreasing derivative is strictly concave by the mean value theorem; conversely, the derivative of a differentiable concave function is nonincreasing, and if it took the same value at two points it would be constant between them, so that the function would be affine on that interval and not strictly concave. This is the equivalence asserted in the statement.

Let $R_\ast$ be strictly concave. For $0<\varepsilon'<\varepsilon_1<\varepsilon_2<\varepsilon_{\max}$, concavity and $R_\ast(\varepsilon')>0$ give
\[
R_\ast(\varepsilon_1)\ \ge\ \frac{\varepsilon_2-\varepsilon_1}{\varepsilon_2-\varepsilon'}\,R_\ast(\varepsilon')+\frac{\varepsilon_1-\varepsilon'}{\varepsilon_2-\varepsilon'}\,R_\ast(\varepsilon_2)\ \ge\ \frac{\varepsilon_1-\varepsilon'}{\varepsilon_2-\varepsilon'}\,R_\ast(\varepsilon_2),
\]
and letting $\varepsilon'\to0$ gives $R_\ast(\varepsilon_1)/\varepsilon_1\ge R_\ast(\varepsilon_2)/\varepsilon_2$: the function $\varepsilon\mapsto R_\ast(\varepsilon)/\varepsilon$ is nonincreasing on $(0,\varepsilon_{\max})$. If it took the same value $c$ at two points $\varepsilon_1<\varepsilon_2$, then, being nonincreasing, it would equal $c$ on $[\varepsilon_1,\varepsilon_2]$; thus $R_\ast(\varepsilon)=c\,\varepsilon$ on this interval, which contradicts strict concavity. Hence $R_\ast(\varepsilon)/\varepsilon$ is strictly decreasing. By (i), $R_\ast(\varepsilon)\to\varepsilon_{\max}$ as $\varepsilon\to\varepsilon_{\max}$, so the ratio tends to $1$ at the right end point. Therefore $R_\ast(\varepsilon)>\varepsilon$ on $(0,\varepsilon_{\max})$, that is, $R_1(u)>u$ on $(0,\xi_{\rm out})$, and again $R_1$ has no fixed point there.

In both cases $\xi_{\rm out}$ is the only crest by (i), so $\Gamma_{\rm out}$ is the only cycle.
\end{proof}
Neither hypothesis of (iii) is proved. Moreover, in the floating-point computations described in Section~\ref{sec:negative}, each of the two hypotheses fails at some parameter pairs with very small or very large $\mu$, while both hold at the remaining pairs examined. Part (iii) is therefore a conditional criterion which cannot be expected to apply on all of $\cP$; a weaker sufficient condition is discussed in Section~\ref{sec:negative}. Consistency tests of the statements of this section are reported in Section~\ref{sec:evidence}.

The next statement records what the planar reduction gives for Dulac functions on $\Dsc$, as opposed to the cycle-average form (A8)(f). For $\Lambda\in C^1(V)$, $V\subset\R^2$ open, put $\Dt\Lambda=\nabla\Lambda\cdot F$, the derivative of $\Lambda$ along the reduced flow, and write $\eta$ for $\eta\circ(\pi|_\cA)^{-1}$ on $\Dsc$; then $\diver(e^{\Lambda}F)=e^{\Lambda}(\Dt\Lambda+\eta-\mu)$ on $V\cap\inter\Dsc$.

\begin{corollary}[Two constraints on Dulac functions]\label{att:cor:dulac}
Assume $k>\kH$, let $V\subset\R^2$ be open and $\Lambda\in C^1(V)$.
\begin{enumerate}[label=\textup{(\roman*)}]
\item If $(0,0)\in V$, then $\Dt\Lambda+\eta-\mu$ equals $\eta_E-\mu=\gamma_E-A>0$ at $(0,0)$ and is positive on a neighbourhood of $(0,0)$ in $\Dsc$. No function that is $C^1$ near the equilibrium satisfies $\Dt\Lambda+\eta-\mu\le0$ near the equilibrium.
\item If $\mathcal J_{\rm in}\subset V$, then every neighbourhood of $\mathcal J_{\rm in}$ contains a point $q\in\operatorname{ins}\mathcal J_{\rm in}$ with $(\Dt\Lambda+\eta-\mu)(q)<0$. This holds whether or not $\Gamma_{\rm in}$ is hyperbolic.
\end{enumerate}
\end{corollary}
\begin{proof}
(i) $F(0,0)=0$, so $\Dt\Lambda(0,0)=0$; $\eta(0,0)=\eta_E$ by (A2), and $\eta_E-\mu=\gamma_E-A>0$ for $k>\kH$. The function $\Dt\Lambda+\eta-\mu$ is continuous on $V\cap\Dsc$ by (A4).

(ii) For $0<u<\xi_{\rm in}$ the point $(u,0)$ lies in $\operatorname{ins}\mathcal J_{\rm in}$ by Lemma~\ref{att:lem:cyclegeom}, and it is not periodic because $\xi_{\rm in}$ is the smallest crest. As in the proof of (A8)(c), its orbit stays in $\operatorname{ins}\mathcal J_{\rm in}$ and has $\omega$-limit set $\mathcal J_{\rm in}$. Its successive crossings of the half-axis $\{v=0,u>0\}$ are the points $(R_1^i(u),0)$, $i\ge0$, and by Lemma~\ref{att:lem:cyclegeom} they satisfy $0<R_1^i(u)<\xi_{\rm in}$. By Lemma~\ref{att:lem:axis}, $F=(0,-g_\ast(s))$ with $g_\ast(s)>0$ at every point $(s,0)$ with $0<s\le\xi_{\rm out}$; hence every closed segment of the $u$-axis contained in $\{(s,0):0<s\le\xi_{\rm out}\}$ is a transversal in the sense of (I2b). By (I2b) the sequence $R_1^i(u)$ is monotone, and it is strictly monotone because the orbit is not periodic. It accumulates at $\xi_{\rm in}$: the orbit accumulates at the point $(\xi_{\rm in},0)$ of its $\omega$-limit set, where $F$ is transverse to the axis, so it crosses the half-axis at points arbitrarily close to $(\xi_{\rm in},0)$. A strictly monotone sequence in $(0,\xi_{\rm in})$ which accumulates at $\xi_{\rm in}$ is increasing. Hence
\begin{equation}\label{att:eq:outward}
R_1(u)>u\qquad(0<u<\xi_{\rm in}).
\end{equation}
Suppose that $V'\subset V$ is a neighbourhood of $\mathcal J_{\rm in}$ with $\Dt\Lambda+\eta-\mu\ge0$ on $V'\cap\operatorname{ins}\mathcal J_{\rm in}$. The functions $R_1$ and $T_1$ are continuous at $\xi_{\rm in}$, with $R_1(\xi_{\rm in})=\xi_{\rm in}$ and $T_1(\xi_{\rm in})$ the period of $\Gamma_{\rm in}$: if $\xi_{\rm in}<\xi_{\rm out}$ this follows from the $C^1$ statement of Proposition~\ref{prop:crestaxis}(i) on $(0,\xi_{\rm out})$, because $\mathcal J_{\rm in}$ meets the half-axis only at its crest and so returns to it for the first time after one period; if $\xi_{\rm in}=\xi_{\rm out}$ it is the statement of Proposition~\ref{prop:crestaxis}(i) at the end point. By continuity of the flow there is $\varepsilon_0>0$ such that, for $\xi_{\rm in}-\varepsilon_0<u<\xi_{\rm in}$, the arc $\{\chi_t(u,0):0\le t\le T_1(u)\}$ lies in $V'\cap\operatorname{ins}\mathcal J_{\rm in}$ and the closed segment $\{(s,0):\xi_{\rm in}-\varepsilon_0\le s\le\xi_{\rm in}\}$ lies in $V'$. For such $u$,
\[
\iota(u)=\int_0^{T_1(u)}(\Dt\Lambda+\eta-\mu)(\chi_t(u,0))\,\mathrm dt-\Lambda(R_1(u),0)+\Lambda(u,0)\ \ge\ \Lambda(u,0)-\Lambda(R_1(u),0),
\]
and \eqref{att:eq:returnderivative} gives $g_\ast(R_1(u))\,e^{\Lambda(R_1(u),0)}R_1'(u)\ge g_\ast(u)\,e^{\Lambda(u,0)}$. With $P_\Lambda(u)=\int_{\xi_{\rm in}-\varepsilon_0}^ug_\ast(s)e^{\Lambda(s,0)}\,\mathrm ds$, which is strictly increasing, this says that $u\mapsto P_\Lambda(R_1(u))-P_\Lambda(u)$ is nondecreasing on $(\xi_{\rm in}-\varepsilon_0,\xi_{\rm in})$; for these $u$ we have $u<R_1(u)<\xi_{\rm in}$ by \eqref{att:eq:outward} and by the bound on the crossings stated before it, so $P_\Lambda$ is defined and differentiable at $R_1(u)$. This function tends to $0$ as $u\uparrow\xi_{\rm in}$. Hence $P_\Lambda(R_1(u))\le P_\Lambda(u)$, that is $R_1(u)\le u$, contradicting \eqref{att:eq:outward}.
\end{proof}
Part (ii) is the return-map form of the Bendixson--Dulac flux argument applied to the region between $\mathcal J_{\rm in}$ and one turn of an orbit spiralling towards it. It does not use hyperbolicity: the cycle-average identity \eqref{att:eq:dulacmean} alone would leave the case $\lambda_c(\Gamma_{\rm in})=0$ open.

\begin{remark}[Divergence in other planes of the attractor]\label{att:rem:planes}
Put $w_1=p+\mu U$ and $w_2=p-(y-k)+\mu U$, so that $w_1'=k-y$ and $w_2'=-k(e^{U}-1)$. On $\inter\Dsc$ the maps $(U,p)\mapsto(U,y)$, $(p,y)$, $(y,w_1)$, $(U,w_2)$, with $y=h(U,p)$, are local $C^1$ diffeomorphisms: their Jacobian determinants are, up to sign, $\eta$, $\zeta$, $\zeta+\mu\eta$ and $1+\eta$. In these planes the reduced field and its divergence are, respectively,
\[
\begin{aligned}
&U'=\mathfrak p_1(U,y),\quad y'=ke^{U}-y;&&\diver=\zeta/\eta-1;\\
&p'=k-\mu p-y,\quad y'=ke^{\mathfrak u_2(p,y)}-y;&&\diver=x/\zeta-A;\\
&y'=ke^{\mathfrak u_3(y,w_1)}-y,\quad w_1'=k-y;&&\diver=x/(\zeta+\mu\eta)-1;\\
&U'=\mathfrak p_4(U,w_2),\quad w_2'=-k(e^{U}-1);&&\diver=(\zeta-\mu)/(1+\eta).
\end{aligned}
\]
Here $\mathfrak p_1,\mathfrak u_2,\mathfrak u_3,\mathfrak p_4$ are defined on $\R^2$ by solving $H(U,p)=y$ for $p$ or for $U$, $H(U,w_1-\mu U)=y$ for $U$, and $p-H(U,p)+k+\mu U=w_2$ for $p$. Each of these equations has exactly one solution for all real values of the other two variables. Indeed, every affine function in the maximum that defines $H$ has $u$-slope $\zeta(\mathbf a)\ge\zeta_->0$ and $v$-slope $-\eta(\mathbf a)\le-\eta_-<0$; comparing $H$ at two points with the affine function that attains the maximum at one of them gives, for $u_1<u_2$ and $v_1<v_2$,
\[
H(u_2,v)-H(u_1,v)\ge\zeta_-(u_2-u_1),\qquad H(u,v_1)-H(u,v_2)\ge\eta_-(v_2-v_1).
\]
Hence the continuous functions $p\mapsto H(U,p)$, $U\mapsto H(U,p)$, $U\mapsto H(U,w_1-\mu U)$ and $p\mapsto p-H(U,p)$ are strictly monotone, and their increments are at least $\eta_-$, $\zeta_-$, $\zeta_-+\mu\eta_-$ and $1+\eta_-$ times the increment of the argument in absolute value; so they tend to $\pm\infty$ at the two ends and each is a bijection of $\R$ onto $\R$. Since $H=h$ is $C^1$ on $\inter\Dsc$, the implicit function theorem shows that the four functions are $C^1$ on the images of $\inter\Dsc$ under the respective maps. The functions $\mathfrak p_1$ and $\mathfrak p_4$ are convex and $\mathfrak u_2$, $\mathfrak u_3$ are concave, since their epigraphs, respectively hypographs, are the sublevel sets $\{H(U,p)-y\le0\}$, $\{H(U,p)-p-k-\mu U+w_2\le0\}$ and $\{H(U,w_1-\mu U)-y\le0\}$ of convex functions. Each divergence differs from $\eta-\mu$ by the flow derivative of the logarithm of the Jacobian:
\begin{gather*}
\Dt\log\eta=\frac\zeta\eta+\mu-1-\eta,\qquad \Dt\log\zeta=\frac x\zeta-1-\eta,\\
\Dt\log(\zeta+\mu\eta)=\frac{x}{\zeta+\mu\eta}-(1+\eta-\mu),\qquad \Dt\log(1+\eta)=\frac{\zeta-\mu}{1+\eta}+\mu-\eta .
\end{gather*}
Hence on every cycle
\[
\lambda_c=\avg\eta-\mu=\avg{\frac\zeta\eta}-1=\avg{\frac x\zeta}-A=\avg{\frac{x}{\zeta+\mu\eta}}-1=\avg{\frac{\zeta-\mu}{1+\eta}} .
\]
\end{remark}
